%% ============================================================================
%% Chapter 13: gRPC & High-Performance RPC
%% Mastering APIs: From Zero to Production Guru
%% ============================================================================
\chapter{gRPC \& High-Performance RPC}
\label{ch:grpc}

%% ----------------------------------------------------------------------------
%% Executive Summary
%% ----------------------------------------------------------------------------
Every API style you have studied so far --- REST over JSON, GraphQL, event-driven
messaging --- ultimately pays the same two taxes on every call: the cost of turning
an in-memory request into bytes, and the cost of moving those bytes across the
network. On public, browser-facing, low-frequency paths those taxes are noise. On
\emph{hot internal paths} --- service-to-service calls inside a data center,
fan-out queries that touch fifty microservices per user request, telemetry streams
emitting thousands of frames per second --- they dominate. gRPC exists to minimize
both taxes simultaneously: it pairs a compact binary serialization format
(Protocol Buffers, Chapter~4) with a multiplexed binary transport (HTTP/2,
Chapter~1) and generates the client and server plumbing from a single
language-neutral interface definition~\cite{grpcdocs,protobufdocs}.

This chapter treats gRPC with the same wire-level rigor we applied to HTTP in
Chapter~1. You will learn the RPC conceptual model --- stubs, skeletons,
marshalling, interface definition --- and precisely why REST+JSON loses on hot
internal paths: payload size, text parsing cost, weak contracts, and the absence
of first-class streaming. You will study gRPC's exact message framing: every
message on an HTTP/2 stream is prefixed with five bytes --- a one-byte compression
flag and a four-byte big-endian length --- followed by the protobuf payload, with
the method identified by \texttt{:path} of the form
\texttt{/package.Service/Method} and the outcome reported in the
\texttt{grpc-status} trailer~\cite{rfc9113,grpcdocs}. You will master Protocol
Buffers as an IDL, all four RPC modes (unary, server streaming, client streaming,
bidirectional streaming), deadline propagation across hops, metadata and
interceptors, the seventeen canonical status codes with rich error details, and
the operational realities that textbooks love to skip: load balancing long-lived
HTTP/2 connections, connection aging, keepalive, server reflection, browser
exposure via grpc-web and JSON transcoding, and the health-checking protocol.

All code is offline-first and deterministic: a synthetic \emph{Northwind
Robotics} telemetry service exercised entirely on localhost and in-process, with
a pytest suite covering deadline exceeded, cancellation, and stream error paths.

\begin{keyconceptbox}
gRPC is not ``REST but faster.'' It is a different \emph{contract discipline}: a
single \texttt{.proto} file is the source of truth from which typed clients and
servers are generated in a dozen languages, the wire format is self-describing by
tag number rather than by field name, and the transport multiplexes many
concurrent streams over one TCP connection. You adopt gRPC to buy contract
safety, streaming semantics, and wire efficiency on internal paths --- and you
pay for it with browser hostility, opaque-on-the-wire debugging, and load
balancing subtleties that REST never had.
\end{keyconceptbox}

%% ----------------------------------------------------------------------------
\section{Learning Objectives \& Prerequisites}
\label{sec:ch13-objectives}

\subsection*{Learning Objectives}
After completing this chapter, its lab, and its exercises, you will be able to:

\begin{enumerate}[label=\textbf{LO-\arabic*.}, leftmargin=2.6cm]
  \item Explain the RPC conceptual model --- client stub, server skeleton,
        marshalling, interface definition language --- and articulate precisely why
        REST+JSON is the wrong tool for hot internal paths (payload size, parse
        cost, codegen, streaming)~\cite{kleppmann2017ddia,newman2021microservices}.
  \item Decode gRPC's wire format by hand: the 5-byte message prefix (1-byte
        compressed flag + 4-byte big-endian length), the
        \texttt{/package.Service/Method} \texttt{:path}, the
        \texttt{application/grpc} content type, trailers-only responses, and the
        \texttt{grpc-status} / \texttt{grpc-message} trailers~\cite{rfc9113,grpcdocs}.
  \item Author Protocol Buffers IDL with messages, services, field rules,
        well-known types, \texttt{oneof}, and \texttt{map}; drive the
        \texttt{protoc} / \texttt{grpcio-tools} codegen pipeline; and apply the
        schema-evolution rules (tag-number discipline, \texttt{reserved}) from
        Chapter~4~\cite{protobufdocs}.
  \item Select and implement the correct RPC mode --- unary, server streaming,
        client streaming, bidirectional streaming --- including the cancellation
        semantics of each.
  \item Engineer deadlines correctly: propagate remaining budget across hops via
        the \texttt{grpc-timeout} header, distinguish per-attempt from overall
        deadlines, and defend servers with \texttt{context.time\_remaining()}
        checks.
  \item Use metadata (headers, trailers, binary \texttt{-bin} keys) and
        client/server interceptors for auth, logging, metrics, and retry policy.
  \item Apply the seventeen canonical gRPC status codes with sound judgment,
        attach rich error details (\texttt{google.rpc.Status}-style payloads in
        binary trailers), and map codes to HTTP statuses when transcoding.
  \item Operate gRPC in production: diagnose the L4 load-balancer hotspot,
        configure client-side and lookaside balancing, tune connection age and
        keepalive, use server reflection with \texttt{grpcurl}, expose services
        to browsers via grpc-web or JSON transcoding, and implement the standard
        health-checking protocol.
\end{enumerate}

\subsection*{Prerequisites}
This chapter assumes you have completed:

\begin{itemize}
  \item \textbf{Chapters 0--2} --- Python 3.11+ environment, virtual environments,
        \texttt{pip}/\texttt{pytest} from PowerShell, and consuming HTTP APIs.
  \item \textbf{Chapter 1} --- the HTTP wire protocol, \emph{especially} the
        HTTP/2 material: streams, frames, HPACK-compressed header blocks,
        \texttt{:path} and other pseudo-headers, flow control, and
        \texttt{RST\_STREAM}. We reuse that vocabulary without re-deriving it.
  \item \textbf{Chapter 3} --- REST's constraints, as the point of contrast.
  \item \textbf{Chapter 4} --- serialization and data contracts: the Protocol
        Buffers wire format (varints, wire types, tag computation), and the
        schema-evolution rules. This chapter assumes them and adds the IDL and
        service layer on top.
\end{itemize}

No network access is required anywhere: every listing binds to
\texttt{localhost} or runs fully in-process, and all data is synthetic, drawn
from the fictional Northwind Robotics warehouse-fleet universe.

\begin{warningbox}
gRPC's performance reputation is real but conditional. The wins come from
binary framing, HPACK header compression, connection reuse, and codegen --- all
of which you can squander with per-call channels, unbounded streams, missing
deadlines, or a naive L4 load balancer. Read Section~\ref{sec:ch13-pitfalls}
before you benchmark.
\end{warningbox}

%% ----------------------------------------------------------------------------
\section{Deep Technical Mechanics \& Architectural Concepts}
\label{sec:ch13-mechanics}

\subsection{The RPC Conceptual Model: Stubs, Skeletons, Marshalling}
\label{subsec:ch13-rpc-model}

\begin{definition}[Remote Procedure Call]
A \emph{remote procedure call} is an invocation that \emph{looks} like a local
function call at the call site but whose execution crosses a process or machine
boundary. The illusion is constructed from four pieces: an \textbf{interface
definition} that names the procedures and their parameter/return types; a
\textbf{client stub} that exposes those signatures locally and converts
arguments to bytes (\emph{marshalling}); a \textbf{server skeleton} (or
\emph{servicer}) that converts bytes back to arguments and invokes the real
implementation; and a \textbf{transport} that moves the bytes.
\end{definition}

The model dates to the early 1980s and has been reincarnated as ONC RPC, DCE/RPC,
CORBA, Java RMI, XML-RPC, SOAP/WSDL, Apache Thrift, and now gRPC. Every
incarnation converges on the same architecture because the problem is invariant:
two processes need a \emph{shared, machine-checkable contract} and a
\emph{mechanical translation} between the contract and each language's native
types. Where the generations differ is in the quality of the contract (IDL
expressiveness and evolution rules), the efficiency of the encoding, and the
semantics of the transport. gRPC's particular bet is: proto3 as the IDL,
protobuf binary as the encoding, HTTP/2 as the transport, and code generation as
the only supported consumption style~\cite{grpcdocs,protobufdocs}.

Two consequences of the stub model deserve emphasis. First, \textbf{the contract
precedes the implementation}: you write \texttt{telemetry.proto} before you
write Python, and both sides generate from it, so client and server cannot
silently disagree about field types. Second, \textbf{the call is not local}: the
stub hides the network but not its physics --- latency, partial failure,
deadlines, and cancellation are first-class in the generated API precisely
because pretending they do not exist is the classic RPC fallacy
\cite{kleppmann2017ddia,newman2021microservices}.

\subsection{Why REST+JSON Loses on Hot Internal Paths}
\label{subsec:ch13-why-not-rest}

REST over JSON is the right default for public APIs --- Chapter~3 argued this at
length --- but on hot internal paths it concedes four compounding
disadvantages:

\begin{enumerate}
  \item \textbf{Payload size.} JSON repeats every field name in every message,
        encodes numbers as decimal text, and base64-inflates binary data. A
        protobuf \texttt{TelemetryFrame} carries field tags as single varint
        bytes (Chapter~4) and is routinely $3$--$10\times$ smaller than the
        equivalent JSON, which translates directly into bandwidth and
        serialization CPU.
  \item \textbf{Parse cost.} JSON parsing is text scanning; protobuf decoding is
        varint reads and memory copies. At tens of thousands of messages per
        second per core, the difference is measurable in tail latency.
  \item \textbf{No enforced contract.} OpenAPI describes REST contracts, but
        nothing forces both sides to generate from it; drift is routine. With
        gRPC the \texttt{.proto} \emph{is} the API: both sides are generated
        artifacts of one file, so drift is a compile error, not an incident.
  \item \textbf{No first-class streaming.} REST has one request and one
        response; streaming is bolted on (SSE, chunked bodies, WebSockets) with
        per-framework semantics. gRPC's four modes make streaming part of the
        contract, with uniform cancellation, flow control, and error semantics.
\end{enumerate}

The honest counterpoint --- developed in Section~\ref{sec:ch13-pitfalls} and the
interview vault --- is that gRPC's advantages are \emph{internal-path}
advantages. Browsers cannot speak raw gRPC, humans cannot read protobuf on the
wire, and the ecosystem of caches, gateways, and tooling built for HTTP/1.1+JSON
mostly does not apply. Choose per path, not per organization.

\subsection{HTTP/2 as the Transport}
\label{subsec:ch13-http2}

Chapter~1 covered HTTP/2's mechanics: one TCP connection carries many
concurrent \emph{streams}; each stream is a sequence of \emph{frames}
(\texttt{HEADERS}, \texttt{DATA}, \texttt{RST\_STREAM}, \texttt{WINDOW\_UPDATE},
\ldots); header blocks are HPACK-compressed; flow control is per-stream and
per-connection via credit windows~\cite{rfc9113}. gRPC maps onto this cleanly:

\begin{center}
\emph{one RPC} $\longleftrightarrow$ \emph{one HTTP/2 stream}.
\end{center}

That single sentence explains most of gRPC's performance profile. Many RPCs
share one TCP connection (no head-of-line blocking between calls, no
connection-pool gymnastics), streams are cheap so unary calls and long-lived
bidi streams coexist, and flow control plus \texttt{RST\_STREAM} give every RPC
independent backpressure and cancellation.

\subsubsection{gRPC Message Framing: The 5-Byte Prefix}
\label{subsubsec:ch13-framing}

HTTP/2 \texttt{DATA} frames are arbitrary byte containers; gRPC imposes its own
message boundary inside them. Each protobuf message --- request or response ---
is serialized on the wire as a \textbf{Length-Prefixed Message}:

\begin{minted}{text}
 byte 0        : Compressed-Flag   (0 = uncompressed, 1 = compressed)
 bytes 1..4    : Message-Length    (uint32, BIG-ENDIAN, network byte order)
 bytes 5..N    : Message           (serialized protobuf, exactly Length bytes)
\end{minted}

For example, a request message that serializes to the 3 protobuf bytes
\texttt{0A 07 ... } of length \texttt{0x0000000B} would appear on the wire as:

\begin{minted}{text}
 00 00 00 00 0B 0A 07 ...   <- flag=0, length=11, then 11 protobuf bytes
\end{minted}

Three properties follow. First, framing is \emph{length-delimited}, so a single
HTTP/2 stream can carry any number of messages in order --- this is what makes
the three streaming modes possible with zero additional machinery. Second, the
compression flag is per-message, so a server may compress opportunistically.
Third, a message may span many HTTP/2 \texttt{DATA} frames, or many messages may
share one frame: the 5-byte prefix, not the frame boundary, is the unit of
meaning.

\subsubsection{Request and Response Anatomy on the Wire}
\label{subsubsec:ch13-anatomy}

A gRPC call is an HTTP/2 \emph{POST} with specific pseudo-headers and a specific
life cycle~\cite{grpcdocs,rfc9113}:

\begin{minted}{text}
 HEADERS (stream 7, flags: END_HEADERS)
   :method    = POST
   :scheme    = https
   :path      = /northwind.telemetry.v1.TelemetryService/GetRobotStatus
   :authority = telemetry.northwind.example
   content-type = application/grpc
   grpc-timeout = 100m
   authorization = Bearer ...
   x-correlation-id = corr-000001
 DATA    (stream 7, flags: END_STREAM)   <- 5-byte prefix + protobuf request
\end{minted}

The \texttt{:path} is exactly \texttt{/}\emph{package.Service}\texttt{/}%
\emph{Method} --- the fully qualified proto name. \texttt{content-type} begins
with \texttt{application/grpc} (commonly \texttt{application/grpc+proto};
\texttt{+json} exists for exotic deployments). The response:

\begin{minted}{text}
 HEADERS (stream 7)               <- "initial metadata"
   :status = 200
   content-type = application/grpc
 DATA    (stream 7)               <- 5-byte prefix + protobuf response(s)
 HEADERS (stream 7, flags: END_STREAM)   <- "trailers"
   grpc-status = 0
   grpc-message = ...
   error-details-bin = <base64>
\end{minted}

Note the asymmetry that surprises every newcomer: \textbf{the HTTP status is
always 200 OK} (barring transport-level failures), and the \emph{real} outcome
of the RPC arrives in the \texttt{grpc-status} \textbf{trailer}. It must be a
trailer because streaming responses may have already emitted dozens of DATA
frames before the failure that determines the final status. When a call fails
before any response message is produced, servers may collapse the whole response
into a single \texttt{HEADERS} frame carrying \texttt{:status 200},
\texttt{grpc-status}, and \texttt{END\_STREAM} --- the \textbf{trailers-only
response}, which is what most error responses look like on the wire.
Figure~\ref{fig:ch13-call-anatomy} shows the complete anatomy.

\begin{figure}[htbp]
\centering
\begin{tikzpicture}[font=\small, >=stealth,
    life/.style={draw=packtgray, dashed},
    box/.style={draw, fill=blue!8, minimum width=3.6cm, inner sep=4pt},
    req/.style={->, thick, packtorange},
    res/.style={->, thick, darkblue},
    lbl/.style={midway, align=left, font=\footnotesize\ttfamily, fill=white, inner sep=1.5pt}]
  \node[box] (client) {Client stub};
  \node[box, right=8.6cm of client] (server) {gRPC server};
  \draw[life] (client.south) -- ++(0,-8.3) coordinate (cbot);
  \draw[life] (server.south) -- ++(0,-8.3) coordinate (sbot);

  \draw[req] ($(client.south)+(0,-0.6)$) -- node[lbl]
      {HEADERS :path=/pkg.Svc/Method, content-type=application/grpc,\\ grpc-timeout=100m, authorization, x-correlation-id}
      ($(server.south)+(0,-0.6)$);
  \draw[req] ($(client.south)+(0,-1.9)$) -- node[lbl]
      {DATA [00 | 00 00 00 0B | protobuf request] END\_STREAM}
      ($(server.south)+(0,-1.9)$);
  \node[note, anchor=west] at ($(server.south)+(0.25,-2.6)$)
      {unmarshal, invoke servicer, marshal};

  \draw[res] ($(server.south)+(0,-3.5)$) -- node[lbl]
      {HEADERS :status=200, content-type=application/grpc}
      ($(client.south)+(0,-3.5)$);
  \draw[res] ($(server.south)+(0,-4.6)$) -- node[lbl]
      {DATA [00 | 00 00 00 21 | protobuf response]}
      ($(client.south)+(0,-4.6)$);
  \draw[res] ($(server.south)+(0,-5.7)$) -- node[lbl]
      {HEADERS grpc-status=0 END\_STREAM (trailers)}
      ($(client.south)+(0,-5.7)$);

  \draw[decorate, decoration={brace, amplitude=5pt}, packtgray]
      ($(sbot)+(0.35,-0.15)$) -- ($(sbot)+(0.35,7.5)$)
      node[note, midway, right=8pt, align=left]
      {one HTTP/2 stream = one RPC\\ (many streams share one\\ TCP connection)};
\end{tikzpicture}
\caption{Anatomy of a unary gRPC call over one HTTP/2 stream. Client frames in
orange, server frames in blue. Every protobuf message is preceded by the 5-byte
prefix (compression flag + big-endian length); the definitive outcome arrives in
the \texttt{grpc-status} trailer.}
\label{fig:ch13-call-anatomy}
\end{figure}
\subsection{Protocol Buffers as the Interface Definition Language}
\label{subsec:ch13-proto-idl}

Chapter~4 covered the protobuf \emph{wire format}; here we cover the protobuf
\emph{language} as the IDL from which gRPC services are
built~\cite{protobufdocs}. The elements you will use daily:

\begin{itemize}
  \item \textbf{Messages and field rules.} In proto3 every field is optional in
        the serialization sense --- absent fields decode to type defaults
        (\texttt{0}, \texttt{""}, \texttt{false}, empty) --- and
        \texttt{repeated} marks a list. proto3 deliberately removed
        \texttt{required}: a required field is a schema-evolution trap, because
        no consumer can ever stop populating it.
  \item \textbf{Tag numbers.} The integer after \texttt{=} is the field's
        identity on the wire (the name is for humans and codegen). Numbers
        1--15 cost one byte of tag; 16--2047 cost two. \emph{Never reuse a
        retired tag number} --- a stale peer will silently reinterpret old bytes
        as the new field. Retire with \texttt{reserved 9, 10;} and
        \texttt{reserved "heading\_deg";} so the compiler enforces the ban
        (Chapter~4's evolution rules, now load-bearing).
  \item \textbf{Well-known types.} \texttt{google.protobuf.Timestamp},
        \texttt{Duration}, \texttt{Struct}, \texttt{Any}, and wrappers like
        \texttt{DoubleValue} ship with protoc. Prefer \texttt{Timestamp} over
        ``epoch seconds in an int64 named \texttt{time}'' --- it is
        self-describing and range-checked.
  \item \textbf{\texttt{oneof}.} At most one branch is set at a time; setting a
        branch clears the others, and presence is preserved on the wire
        (\texttt{WhichOneof("action")} in Python). Use it for genuine
        alternatives --- our \texttt{location} and \texttt{action} fields ---
        not as a general union type.
  \item \textbf{\texttt{map<K,V>}.} Sugar for a repeated entry message; keys may
        be integral or string types. Map ordering is \emph{not} guaranteed on
        the wire --- never key application logic on it.
  \item \textbf{Services.} A \texttt{service} block names RPCs and their
        cardinality: \texttt{rpc M(Req) returns (Resp)} for unary, with
        \texttt{stream} on either side for the streaming modes. This is the
        only part of the file protoc's gRPC plugin consumes.
\end{itemize}

The codegen pipeline is mechanical and hermetic:

\begin{minted}{text}
python -m grpc_tools.protoc -I. --python_out=. --grpc_python_out=. telemetry.proto
\end{minted}

\texttt{grpc\_tools.protoc} bundles the protoc compiler and the well-known
proto files (so the \texttt{google/protobuf/timestamp.proto} import resolves
offline). It emits two modules: \texttt{telemetry\_pb2.py} (message classes:
builders, accessors, \texttt{SerializeToString}/\texttt{ParseFromString},
\texttt{WhichOneof}, map containers) and \texttt{telemetry\_pb2\_grpc.py}
(the \texttt{TelemetryServiceStub} client class, the
\texttt{TelemetryServiceServicer} base class you subclass, and the
\texttt{add\_TelemetryServiceServicer\_to\_server} registration function).
\textbf{Generated files are build artifacts}: regenerate, never hand-edit, and
regenerate in CI so stubs can never drift from the \texttt{.proto}.

\begin{examplebox}
Schema evolution, one paragraph: you may add fields with \emph{fresh} tag
numbers (old peers ignore unknown tags; new peers see defaults), remove fields
only by reserving their tags and names first, widen \texttt{int32} to
\texttt{int64} (same wire type), and add \texttt{oneof} branches. You may never
reuse tag numbers, change a field's wire type, or move an existing field into a
\texttt{oneof}. Our \texttt{TelemetryFrame} reserves tags 9--10 and the names
\texttt{heading\_deg}/\texttt{speed\_mps}, retired in v1.1 --- the reservation
makes a v1.2 reintroduction a \emph{compile error} instead of a production data
corruption (Chapter~4, Section on tag discipline)~\cite{protobufdocs}.
\end{examplebox}

\subsection{The Four RPC Modes}
\label{subsec:ch13-modes}

Cardinality is declared in the IDL and fixed at compile time; each mode has a
canonical use case and its own cancellation semantics
(Figure~\ref{fig:ch13-four-modes}):

\begin{enumerate}
  \item \textbf{Unary} --- one request, one response. The default for
        command/query shapes (\texttt{GetRobotStatus}). Cancellation: client
        abort sends \texttt{RST\_STREAM}; server observes
        \texttt{context.is\_active() == False}.
  \item \textbf{Server streaming} --- one request, many responses
        (\texttt{StreamTelemetry}: ``subscribe to this robot's frames'').
        Canonical for change feeds, tailing, and large result sets consumed
        incrementally. Cancellation matters most here: an unbounded stream the
        client abandoned is a server-side resource leak unless the handler
        checks \texttt{is\_active()} between yields.
  \item \textbf{Client streaming} --- many requests, one response
        (\texttt{UploadTelemetry}: bulk ingestion). The server reads until
        end-of-stream, then computes one summary. Errors discovered mid-stream
        should abort with status and details, not be silently tallied.
  \item \textbf{Bidirectional streaming} --- two independent message sequences
        on one stream (\texttt{ControlLoop}: commands down, acknowledgements
        up). The sequences are \emph{not} request/response pairs unless your
        protocol makes them so; ordering is guaranteed within each direction
        only. This is the mode for control planes, chat-like protocols, and
        interactive sessions.
\end{enumerate}

\begin{figure}[htbp]
\centering
\begin{tikzpicture}[font=\footnotesize, >=stealth,
    life/.style={draw=packtgray, dashed},
    peer/.style={draw, fill=blue!8, inner sep=2pt, font=\footnotesize\bfseries},
    req/.style={->, thick, packtorange},
    res/.style={->, thick, darkblue},
    cap/.style={font=\footnotesize\bfseries}]

  % ---- Unary (top-left) ----
  \begin{scope}
    \node[modcap] at (1.5,0.3) {Unary};
    \node[peer] (c1) at (0,0) {C};
    \node[peer] (s1) at (3,0) {S};
    \draw[life] (c1.south) -- ++(0,-2.4);
    \draw[life] (s1.south) -- ++(0,-2.4);
    \draw[req] ($(c1.south)+(0,-0.5)$) -- node[above]{1 req} ($(s1.south)+(0,-0.5)$);
    \draw[res] ($(s1.south)+(0,-1.7)$) -- node[above]{1 resp} ($(c1.south)+(0,-1.7)$);
  \end{scope}

  % ---- Server streaming (top-right) ----
  \begin{scope}[xshift=7.2cm]
    \node[modcap] at (1.5,0.3) {Server streaming};
    \node[peer] (c2) at (0,0) {C};
    \node[peer] (s2) at (3,0) {S};
    \draw[life] (c2.south) -- ++(0,-2.4);
    \draw[life] (s2.south) -- ++(0,-2.4);
    \draw[req] ($(c2.south)+(0,-0.5)$) -- node[above]{1 req} ($(s2.south)+(0,-0.5)$);
    \draw[res] ($(s2.south)+(0,-1.1)$) -- ($(c2.south)+(0,-1.1)$);
    \draw[res] ($(s2.south)+(0,-1.6)$) -- node[below]{n resp} ($(c2.south)+(0,-1.6)$);
    \draw[res] ($(s2.south)+(0,-2.1)$) -- ($(c2.south)+(0,-2.1)$);
  \end{scope}

  % ---- Client streaming (bottom-left) ----
  \begin{scope}[yshift=-3.6cm]
    \node[modcap] at (1.5,0.3) {Client streaming};
    \node[peer] (c3) at (0,0) {C};
    \node[peer] (s3) at (3,0) {S};
    \draw[life] (c3.south) -- ++(0,-2.4);
    \draw[life] (s3.south) -- ++(0,-2.4);
    \draw[req] ($(c3.south)+(0,-0.6)$) -- ($(s3.south)+(0,-0.6)$);
    \draw[req] ($(c3.south)+(0,-1.1)$) -- node[below]{n req} ($(s3.south)+(0,-1.1)$);
    \draw[req] ($(c3.south)+(0,-1.6)$) -- ($(s3.south)+(0,-1.6)$);
    \draw[res] ($(s3.south)+(0,-2.2)$) -- node[above]{1 resp} ($(c3.south)+(0,-2.2)$);
  \end{scope}

  % ---- Bidirectional (bottom-right) ----
  \begin{scope}[xshift=7.2cm, yshift=-3.6cm]
    \node[modcap] at (1.5,0.3) {Bidirectional};
    \node[peer] (c4) at (0,0) {C};
    \node[peer] (s4) at (3,0) {S};
    \draw[life] (c4.south) -- ++(0,-2.4);
    \draw[life] (s4.south) -- ++(0,-2.4);
    \draw[req] ($(c4.south)+(0,-0.5)$) -- ($(s4.south)+(0,-0.5)$);
    \draw[res] ($(s4.south)+(0,-0.9)$) -- ($(c4.south)+(0,-0.9)$);
    \draw[req] ($(c4.south)+(0,-1.4)$) -- ($(s4.south)+(0,-1.4)$);
    \draw[res] ($(s4.south)+(0,-1.8)$) -- ($(c4.south)+(0,-1.8)$);
    \node[note, anchor=west] at ($(c4.south)+(0.1,-2.3)$) {interleaved, in-order per direction};
  \end{scope}
\end{tikzpicture}
\caption{Message timelines for the four RPC modes. Client-to-server messages in
orange, server-to-client in blue. Every arrow is a length-prefixed message on
the same HTTP/2 stream; the stream closes with the \texttt{grpc-status}
trailer.}
\label{fig:ch13-four-modes}
\end{figure}

\subsection{Deadlines, Cancellation, and Metadata}
\label{subsec:ch13-deadlines}

\textbf{Deadline mechanics.} When a Python client passes
\texttt{timeout=0.1}, the gRPC core computes an absolute deadline
(\emph{now} + 100 ms) and transmits the \emph{remaining} budget in the
\texttt{grpc-timeout} header as an integer with a unit suffix ---
\texttt{H} hours, \texttt{M} minutes, \texttt{S} seconds, \texttt{m}
milliseconds, \texttt{u} microseconds, \texttt{n} nanoseconds, at most eight
digits (e.g. \texttt{grpc-timeout: 99950u} after 50 $\mu$s of transit). The
server's \texttt{context.time\_remaining()} is derived from that header minus
observed transit time --- so \textbf{the budget propagates}: when your handler
fans out to a downstream gRPC service, forwarding \texttt{time\_remaining()} as
the downstream timeout means the whole call tree shares one end-to-end budget
instead of each hop inventing its own. This is the single most important
resilience behavior in the chapter; the lab exercises it.

\textbf{Per-attempt vs.\ overall.} Retries (Section~\ref{subsec:ch13-code-client})
spend \emph{attempts} from the \emph{same} overall deadline; the service config's
\texttt{timeout} field caps each RPC, while per-attempt budgets are bounded by
backoff configuration. A call that needs 3 attempts $\times$ 300 ms must be
given an overall deadline well above 900 ms, or the retries are decoration.

\textbf{Cancellation.} Client cancellation emits \texttt{RST\_STREAM} (Chapter
1); the server must poll \texttt{context.is\_active()} in loops and generators
and stop producing. Cancellation is advisory and cooperative --- a server that
never checks will happily compute into the void.

\textbf{Metadata.} Request headers and response trailers are both
\emph{metadata}: lowercase ASCII keys, string values --- \emph{except} keys
suffixed \texttt{-bin}, which carry raw binary (base64-encoded on the HTTP/2
wire, exposed as bytes in the API). Auth rides in metadata ---
\texttt{authorization: Bearer \ldots} on every call, attached by an interceptor,
never hand-threaded through business logic --- and rich error details ride in
\texttt{-bin} trailers (Section~\ref{subsec:ch13-errors}). Metadata keys
beginning with \texttt{grpc-} are reserved for the framework.

\subsection{Interceptors: The Cross-Cutting Layer}
\label{subsec:ch13-interceptors}

Interceptors are gRPC's middleware: functions that wrap every call on a channel
(client side) or every handler on a server, for concerns that must be uniform
--- authentication, correlation IDs, logging, metrics, retry hedging, error
mapping. They come in client/server $\times$ unary/stream varieties; Python
models the client ones as the four abstract classes
\texttt{grpc.UnaryUnaryClientInterceptor} \ldots{}
\texttt{grpc.StreamStreamClientInterceptor}, composed with
\texttt{grpc.intercept\_channel}, and the server side as
\texttt{grpc.ServerInterceptor.intercept\_service}, which receives the
not-yet-bound handler and returns a (wrapped) handler.
Figure~\ref{fig:ch13-interceptors} shows the chain; Section~\ref{sec:ch13-code}
implements both ends.

\begin{figure}[htbp]
\centering
\begin{tikzpicture}[font=\footnotesize, >=stealth,
    node distance=0.45cm,
    box/.style={draw, rounded corners=2pt, fill=blue!8, minimum height=0.85cm, inner sep=4pt},
    ibox/.style={box, fill=packtorange!12},
    tport/.style={box, fill=green!10, minimum width=4.4cm},
    flow/.style={->, thick, packtgray}]
  % outbound chain
  \node[box] (app) {Client app};
  \node[ibox, right=of app] (c1) {auth + correlation ID};
  \node[ibox, right=of c1] (c2) {logging / metrics};
  \node[ibox, right=of c2] (c3) {retry policy};
  \node[tport, right=of c3] (chan) {channel (HTTP/2)};
  \draw[flow] (app) -- (c1);
  \draw[flow] (c1) -- (c2);
  \draw[flow] (c2) -- (c3);
  \draw[flow] (c3) -- (chan);

  % inbound chain
  \node[tport, below=1.1cm of chan] (lst) {listener (HTTP/2)};
  \node[ibox, left=of lst] (s1) {error mapping};
  \node[ibox, left=of s1] (s2) {authz check};
  \node[ibox, left=of s2] (s3) {logging / metrics};
  \node[box, left=of s3] (svc) {Service impl};
  \draw[flow] (lst) -- (s1);
  \draw[flow] (s1) -- (s2);
  \draw[flow] (s2) -- (s3);
  \draw[flow] (s3) -- (svc);

  \draw[flow] (chan) -- node[right, note]{one TCP connection, many streams} (lst);
  \draw[flow, dashed] (svc.north) .. controls +(0,0.7) and +(-2.4,0.4) ..
      node[above, note]{response path unwinds the chain in reverse} (app.north -| c1);
\end{tikzpicture}
\caption{Interceptor chain. Client interceptors wrap outbound calls (attach
metadata, observe latency, drive retries); server interceptors wrap inbound
handlers (enforce authz, map domain errors to status codes, emit metrics).}
\label{fig:ch13-interceptors}
\end{figure}

\subsection{The Error Model}
\label{subsec:ch13-errors}

gRPC replaces HTTP's flat status taxonomy with seventeen canonical codes carried
in the \texttt{grpc-status} trailer~\cite{grpcdocs}. The full set, with usage
guidance:

\begin{table}[htbp]
  \centering
  \small
  \begin{tabularx}{\textwidth}{c l X c}
    \toprule
    \textbf{\#} & \textbf{Code} & \textbf{Use when} & \textbf{HTTP map} \\
    \midrule
    0  & OK & Success (not an error). & 200 \\
    1  & CANCELLED & Caller cancelled (or deadline-adjacent abort). & 499 \\
    2  & UNKNOWN & Error from another namespace / unclassified. & 500 \\
    3  & INVALID\_ARGUMENT & Client sent malformed input; retry unchanged will fail. & 400 \\
    4  & DEADLINE\_EXCEEDED & Budget expired (note: server-set, vs.\ CANCELLED for client abort). & 504 \\
    5  & NOT\_FOUND & Named resource does not exist. & 404 \\
    6  & ALREADY\_EXISTS & Create collided with existing resource. & 409 \\
    7  & PERMISSION\_DENIED & Authenticated, but not allowed. & 403 \\
    8  & RESOURCE\_EXHAUSTED & Quota / rate limit / disk full. & 429 \\
    9  & FAILED\_PRECONDITION & System state forbids the op (e.g.\ non-empty dir delete). & 400 \\
    10 & ABORTED & Concurrency conflict; safe to retry whole op (transactions). & 409 \\
    11 & OUT\_OF\_RANGE & Value outside valid range (seek past end). & 400 \\
    12 & UNIMPLEMENTED & Method unknown to this server. & 501 \\
    13 & INTERNAL & Server bug / invariant violation. & 500 \\
    14 & UNAVAILABLE & Transient: server down, overloaded, restarting. \emph{Retryable.} & 503 \\
    15 & DATA\_LOSS & Unrecoverable data corruption. & 500 \\
    16 & UNAUTHENTICATED & Missing/invalid credentials. & 401 \\
    \bottomrule
  \end{tabularx}
  \caption{The 17 canonical gRPC status codes and their conventional HTTP
  mapping when transcoding to JSON APIs (the mapping used by Google's
  \texttt{google/rpc/code.proto} conventions).}
  \label{tab:ch13-status}
\end{table}

Three judgment calls trip up even senior engineers. \textbf{INVALID\_ARGUMENT
vs.\ FAILED\_PRECONDITION vs.\ OUT\_OF\_RANGE}: the first means the request is
malformed no matter the system state; the second means the request is fine but
the system's state forbids it; the third is for values outside a defined range
(e.g.\ seeking past the end of a stream). \textbf{UNAUTHENTICATED
vs.\ PERMISSION\_DENIED}: the first means ``we don't know who you are''
(credentials missing/invalid --- re-auth may fix it); the second means ``we know
who you are and the answer is no'' (re-auth will not help).
\textbf{NOT\_FOUND vs.\ PERMISSION\_DENIED}: when revealing existence would
leak information, return PERMISSION\_DENIED (or even NOT\_FOUND) uniformly ---
but pick one policy per resource class and document it.

\textbf{Rich error details.} A status code plus a string is not an error
\emph{model}. Google APIs standardize a richer payload: a
\texttt{google.rpc.Status} message --- code, message, and \texttt{repeated
google.protobuf.Any details} carrying typed payloads such as
\texttt{BadRequest} (per-field violations), \texttt{QuotaFailure},
\texttt{RetryInfo}, and \texttt{ErrorInfo} (stable machine-readable
\texttt{reason} + \texttt{domain}) --- serialized into the binary trailer
\texttt{grpc-status-details-bin}. Our listing implements the same pattern with a
first-class \texttt{ErrorDetails} message in the chapter proto and an
\texttt{error-details-bin} trailer, so it runs with only
\texttt{grpcio}/\texttt{protobuf} installed; swapping in
\texttt{googleapis-common-protos} and the standard key is a drop-in upgrade.

\textbf{Transcoding.} When a gateway translates gRPC to HTTP+JSON
(Section~\ref{subsec:ch13-operational}), each status code maps to one HTTP
status (Table~\ref{tab:ch13-status}); details typically serialize into the
response body as a JSON \texttt{error} object. Design your codes and details so
both audiences --- native gRPC clients and transcoded REST clients --- can act
on them.

\subsection{Operational gRPC: Load Balancing, Connections, Tooling}
\label{subsec:ch13-operational}

\textbf{The L4 problem.} gRPC multiplexes \emph{many RPCs over one long-lived
TCP connection}. A Layer-4 (TCP) load balancer balances \emph{connections}, so
once a client connects to a pod, every RPC it ever makes lands on that pod ---
autoscaling adds idle pods while one pod melts (the war story in
Section~\ref{sec:ch13-lab} tells it with numbers). The standard remedies:
\textbf{(a) client-side LB} --- the client resolves all backends (DNS) and
load-balances per-RPC with \texttt{round\_robin} (or per-connection with the
default \texttt{pick\_first}); \textbf{(b) lookaside LB} --- a control plane
(the historical \texttt{grpclb} protocol) hands the client a curated server
list, keeping data-plane traffic direct; \textbf{(c) xDS} --- the CNCF/Envoy
service-discovery API family (LDS/CDS/RDS/EDS) that modern meshes and gRPC's
``proxyless service mesh'' mode use to push routing, LB policy, and config to
clients; and \textbf{(d) L7 proxies} (Envoy, Linkerd) that terminate HTTP/2 and
rebalance per-stream~\cite{grpcdocs,newman2021microservices}.

\textbf{Connection management.} Two server options prevent the pathologies of
immortal connections: \texttt{grpc.max\_connection\_age\_ms} forces periodic
reconnects (clients transparently redial, landing on a rebalanced backend ---
set \texttt{max\_connection\_age\_grace\_ms} to drain in-flight RPCs), and
\texttt{grpc.keepalive\_time\_ms} sends HTTP/2 PINGs so dead peers and idle
connections through NATs/firewalls are detected. Client-side
\texttt{grpc.keepalive\_time\_ms} must exceed the server's
\texttt{min\_ping\_interval} policy or the server will throttle you with
GOAWAYs.

\textbf{Server reflection and grpcurl.} A server opting into the reflection
service can enumerate its services and descriptors at runtime, which powers
\texttt{grpcurl} --- ``curl for gRPC'' --- with no local \texttt{.proto} copy.
In Python, enable it with the \texttt{grpcio-reflection} package; the lab shows
the exact invocations.

\textbf{Browsers: grpc-web and JSON transcoding.} Browsers cannot open raw
HTTP/2 streams with trailer access, so two bridges exist: \textbf{grpc-web}
speaks a framed variant of the protocol through a proxy (Envoy has a native
filter), giving web clients generated stubs and unary/server-streaming calls;
\textbf{HTTP/JSON transcoding} annotates the proto
(\texttt{google.api.http}) so a gateway (Envoy's gRPC-JSON transcoder,
gRPC-Gateway) exposes each RPC as a conventional REST endpoint --- the standard
way to keep one internal gRPC contract while publishing a public JSON API.

\textbf{Health checking.} The \texttt{grpc.health.v1.Health} service
(\texttt{Check}, \texttt{Watch}) is the standard liveness/readiness contract:
orchestrators and LBs query it per service name, and servers flip
\texttt{SERVING}/\texttt{NOT\_SERVING} during drains. Implement it from day one
--- a load balancer that cannot distinguish ``up'' from ``serving'' will route
traffic into your deploys.

%% ----------------------------------------------------------------------------
\section{Production-Ready Code Implementation}
\label{sec:ch13-code}
The running example is the \emph{Northwind Robotics Telemetry Service}: a fleet
of warehouse robots emits telemetry frames, operators query status, and a
control loop steers individual units. All six listings below were executed
against Python 3.11+ with \texttt{grpcio}, \texttt{grpcio-tools},
\texttt{protobuf}, and \texttt{pytest} (Appendix, Section~%
\ref{sec:ch13-deps}), entirely offline on \texttt{localhost}. Data is synthetic
and deterministic: frames derive from a fixed seed and a fixed base epoch,
never from the wall clock.

\subsection{Listing 1 --- Interface Definition: \texttt{telemetry.proto}}
\label{subsec:ch13-code-proto}

One file is the whole contract: messages with well-known types
(\texttt{Timestamp}), a \texttt{map}, two \texttt{oneof}s, \texttt{reserved}
tags for retired fields, and a service declaring all four RPC modes.

\begin{minted}{text}
syntax = "proto3";

package northwind.telemetry.v1;

import "google/protobuf/timestamp.proto";

// Operating mode reported by a Northwind Robotics warehouse robot.
enum RobotMode {
  ROBOT_MODE_UNSPECIFIED = 0;
  ROBOT_MODE_IDLE = 1;
  ROBOT_MODE_MOVING = 2;
  ROBOT_MODE_CHARGING = 3;
  ROBOT_MODE_FAULT = 4;
}

message AislePosition {
  string aisle_id = 1;
  double offset_m = 2;
}

message GpsFix {
  double latitude = 1;
  double longitude = 2;
  double hdop = 3;
}

// One telemetry sample emitted by a robot.
message TelemetryFrame {
  string robot_id = 1;
  google.protobuf.Timestamp observed_at = 2;  // well-known type
  double battery_pct = 3;                     // 0.0 .. 100.0
  double chassis_temp_celsius = 4;
  map<string, double> joint_positions_rad = 5;
  RobotMode mode = 6;

  oneof location {                            // exactly one branch is set
    AislePosition aisle = 7;
    GpsFix gps = 8;
  }

  // Tag numbers and names retired in v1.1; never reuse them (see Chapter 4).
  reserved 9, 10;
  reserved "heading_deg", "speed_mps";
}

message RobotStatusRequest {
  string robot_id = 1;
}

message RobotStatusReply {
  string robot_id = 1;
  RobotMode mode = 2;
  double battery_pct = 3;
  google.protobuf.Timestamp firmware_built_at = 4;
  TelemetryFrame last_frame = 5;
}

message StreamTelemetryRequest {
  string robot_id = 1;
  uint32 max_frames = 2;  // hard bound; the server never exceeds it
  uint32 period_ms = 3;   // pacing between frames
}

message UploadSummary {
  uint32 accepted = 1;
  uint32 rejected = 2;
  double min_battery_pct = 3;
  double max_battery_pct = 4;
}

message ControlCommand {
  string robot_id = 1;
  uint64 sequence = 2;

  oneof action {
    RobotMode set_mode = 3;
    bool estop = 4;
    double set_speed_mps = 5;
  }
}

// Structured error payload, carried in the "error-details-bin" trailer.
// Mirrors the shape of google.rpc.Status details used by Google APIs.
message ErrorDetails {
  string reason = 1;
  string domain = 2;  // e.g. "telemetry.northwind.example"
  repeated FieldViolation violations = 3;
  string correlation_id = 4;
}

message FieldViolation {
  string field = 1;
  string description = 2;
}

// The four RPC modes, one of each kind.
service TelemetryService {
  rpc GetRobotStatus(RobotStatusRequest) returns (RobotStatusReply);  // unary
  rpc StreamTelemetry(StreamTelemetryRequest) returns (stream TelemetryFrame);
  rpc UploadTelemetry(stream TelemetryFrame) returns (UploadSummary);
  rpc ControlLoop(stream ControlCommand) returns (stream TelemetryFrame);
}
\end{minted}

Generate the stubs from PowerShell (the bundled \texttt{grpc\_tools.protoc}
resolves the well-known-type import offline):

\begin{minted}{text}
python -m grpc_tools.protoc -I. --python_out=. --grpc_python_out=. telemetry.proto
\end{minted}

This emits \texttt{telemetry\_pb2.py} (messages) and
\texttt{telemetry\_pb2\_grpc.py} (stub, servicer base, registration). The
generated-stub usage pattern is always the same three moves:
\emph{subclass} the servicer and implement the RPCs (Listing~2);
\emph{register} the servicer on a \texttt{grpc.server}; on the client,
\emph{instantiate the stub} against a channel and call methods as if local
(Listing~3).

\subsection{Listing 2 --- Server: All Four Modes with Deadline and Cancellation Discipline}
\label{subsec:ch13-code-server}

Note the defensive habits: every handler checks \texttt{context.is\_active()}
in its loop, re-checks \texttt{context.time\_remaining()} after any stall, and
raises typed \texttt{DomainError}s that Listing~4 maps to status codes with
rich details. \texttt{RB-SLOW} simulates a slow dependency so the deadline test
has something to bite.

\begin{minted}{python}
"""Northwind Robotics telemetry service -- gRPC server implementation.

Chapter 13 listing. Offline-first: binds to localhost only; every payload is
synthetic and derived from a fixed seed, never from the wall clock.
"""
from __future__ import annotations

import argparse
import random
import time
from concurrent import futures
from typing import Iterable, Iterator

import grpc

import telemetry_pb2
import telemetry_pb2_grpc

BASE_EPOCH_S = 1_700_000_000  # fixed synthetic epoch for deterministic payloads
KNOWN_ROBOTS = ("RB-1001", "RB-1002", "RB-1003", "RB-SLOW")
MIN_BUDGET_S = 0.005


class DomainError(Exception):
    """Application error carrying its intended gRPC status mapping."""

    status_code: grpc.StatusCode = grpc.StatusCode.INTERNAL
    reason: str = "INTERNAL"

    def __init__(
        self,
        message: str,
        violations: list[telemetry_pb2.FieldViolation] | None = None,
    ) -> None:
        super().__init__(message)
        self.message = message
        self.violations = violations or []


class UnknownRobotError(DomainError):
    status_code = grpc.StatusCode.NOT_FOUND
    reason = "ROBOT_NOT_FOUND"


class FrameValidationError(DomainError):
    status_code = grpc.StatusCode.INVALID_ARGUMENT
    reason = "FRAME_VALIDATION_FAILED"


def _check_known(robot_id: str) -> None:
    if robot_id not in KNOWN_ROBOTS:
        raise UnknownRobotError(f"unknown robot_id {robot_id!r}")


def _require_budget(
    context: grpc.ServicerContext, min_s: float = MIN_BUDGET_S
) -> None:
    """Fail fast if the propagated deadline leaves no useful budget."""
    remaining = context.time_remaining()
    if remaining < min_s:
        context.abort(
            grpc.StatusCode.DEADLINE_EXCEEDED,
            f"budget exhausted: {remaining * 1000:.2f} ms remaining",
        )


class TelemetryService(telemetry_pb2_grpc.TelemetryServiceServicer):
    """All four RPC modes with explicit deadline and cancellation checks."""

    def __init__(self, seed: int = 1337) -> None:
        self._seed = seed

    def _frame(
        self,
        robot_id: str,
        seq: int,
        mode: telemetry_pb2.RobotMode | None = None,
    ) -> telemetry_pb2.TelemetryFrame:
        # Per-frame seeded RNG keeps payloads deterministic and reproducible.
        rng = random.Random(f"{self._seed}:{robot_id}:{seq}")
        frame = telemetry_pb2.TelemetryFrame(
            robot_id=robot_id,
            battery_pct=round(95.0 - 0.01 * seq + rng.uniform(-0.25, 0.25), 3),
            chassis_temp_celsius=round(32.0 + rng.uniform(-1.0, 1.0), 2),
            mode=(
                telemetry_pb2.ROBOT_MODE_MOVING if mode is None else mode
            ),
        )
        frame.observed_at.FromSeconds(BASE_EPOCH_S + seq * 5)
        frame.joint_positions_rad["shoulder"] = round(rng.uniform(-1.0, 1.0), 4)
        frame.joint_positions_rad["elbow"] = round(rng.uniform(0.0, 2.0), 4)
        frame.aisle.CopyFrom(
            telemetry_pb2.AislePosition(
                aisle_id="A-07", offset_m=round(rng.uniform(0.0, 120.0), 2)
            )
        )
        return frame

    # -- Mode 1: unary ----------------------------------------------------
    def GetRobotStatus(
        self,
        request: telemetry_pb2.RobotStatusRequest,
        context: grpc.ServicerContext,
    ) -> telemetry_pb2.RobotStatusReply:
        _require_budget(context)
        _check_known(request.robot_id)
        if request.robot_id == "RB-SLOW":
            time.sleep(0.200)  # simulated slow dependency
            _require_budget(context, min_s=0.100)  # re-check after the stall
        last = self._frame(request.robot_id, seq=0)
        reply = telemetry_pb2.RobotStatusReply(
            robot_id=request.robot_id,
            mode=last.mode,
            battery_pct=last.battery_pct,
        )
        reply.firmware_built_at.FromSeconds(BASE_EPOCH_S - 86_400)
        reply.last_frame.CopyFrom(last)
        return reply

    # -- Mode 2: server streaming ------------------------------------------
    def StreamTelemetry(
        self,
        request: telemetry_pb2.StreamTelemetryRequest,
        context: grpc.ServicerContext,
    ) -> Iterator[telemetry_pb2.TelemetryFrame]:
        _check_known(request.robot_id)
        if request.max_frames == 0:
            raise FrameValidationError(
                "max_frames must be >= 1",
                [
                    telemetry_pb2.FieldViolation(
                        field="max_frames", description="must be >= 1"
                    )
                ],
            )
        period_s = max(request.period_ms, 1) / 1000.0
        for seq in range(request.max_frames):
            if not context.is_active():
                return  # client cancelled: stop producing promptly
            _require_budget(context)
            time.sleep(period_s)
            yield self._frame(request.robot_id, seq)

    # -- Mode 3: client streaming -------------------------------------------
    def UploadTelemetry(
        self,
        request_iterator: Iterable[telemetry_pb2.TelemetryFrame],
        context: grpc.ServicerContext,
    ) -> telemetry_pb2.UploadSummary:
        accepted = 0
        rejected = 0
        min_battery = float("inf")
        max_battery = float("-inf")
        for frame in request_iterator:
            if not context.is_active():
                break
            if frame.robot_id not in KNOWN_ROBOTS:
                raise FrameValidationError(
                    "frame references unknown robot",
                    [
                        telemetry_pb2.FieldViolation(
                            field="robot_id",
                            description="must be a known robot",
                        )
                    ],
                )
            if not 0.0 <= frame.battery_pct <= 100.0:
                rejected += 1
                continue
            accepted += 1
            min_battery = min(min_battery, frame.battery_pct)
            max_battery = max(max_battery, frame.battery_pct)
        return telemetry_pb2.UploadSummary(
            accepted=accepted,
            rejected=rejected,
            min_battery_pct=round(min_battery, 3) if accepted else 0.0,
            max_battery_pct=round(max_battery, 3) if accepted else 0.0,
        )

    # -- Mode 4: bidirectional streaming -------------------------------------
    def ControlLoop(
        self,
        request_iterator: Iterable[telemetry_pb2.ControlCommand],
        context: grpc.ServicerContext,
    ) -> Iterator[telemetry_pb2.TelemetryFrame]:
        for command in request_iterator:
            if not context.is_active():
                return
            _require_budget(context)
            _check_known(command.robot_id)
            action = command.WhichOneof("action")
            if action == "set_mode":
                mode = command.set_mode
            elif action == "estop":
                mode = (
                    telemetry_pb2.ROBOT_MODE_FAULT
                    if command.estop
                    else telemetry_pb2.ROBOT_MODE_IDLE
                )
            elif action == "set_speed_mps":
                mode = telemetry_pb2.ROBOT_MODE_MOVING
            else:
                raise FrameValidationError(
                    "command has no action set",
                    [
                        telemetry_pb2.FieldViolation(
                            field="action", description="oneof must be set"
                        )
                    ],
                )
            yield self._frame(command.robot_id, seq=command.sequence, mode=mode)


def make_server(
    interceptors: tuple[grpc.ServerInterceptor, ...] = (),
    max_workers: int = 8,
) -> grpc.Server:
    server = grpc.server(
        futures.ThreadPoolExecutor(max_workers=max_workers),
        interceptors=interceptors,
        options=[
            ("grpc.max_connection_age_ms", 300_000),  # let L4/L7 LBs rebalance
            ("grpc.keepalive_time_ms", 60_000),
            ("grpc.keepalive_permit_without_calls", 0),
        ],
    )
    telemetry_pb2_grpc.add_TelemetryServiceServicer_to_server(
        TelemetryService(), server
    )
    return server


def main() -> None:
    from error_middleware import ErrorMappingInterceptor

    parser = argparse.ArgumentParser(description="Telemetry gRPC server")
    parser.add_argument("--port", type=int, default=50051)
    args = parser.parse_args()
    server = make_server(interceptors=(ErrorMappingInterceptor(),))
    bound = server.add_insecure_port(f"localhost:{args.port}")
    server.start()
    print(f"telemetry server listening on localhost:{bound}")
    server.wait_for_termination()


if __name__ == "__main__":
    main()
\end{minted}

\subsection{Listing 3 --- Error-Mapping Middleware with Rich Details}
\label{subsec:ch13-code-middleware}

A \texttt{grpc.ServerInterceptor} rebuilds each handler so that any
\texttt{DomainError} becomes the right status code plus a serialized
\texttt{ErrorDetails} message in the binary \texttt{error-details-bin} trailer
--- the same transport trick Google APIs use for \texttt{google.rpc.Status} in
\texttt{grpc-status-details-bin}. \texttt{decode\_error\_details} is the client
half; the test suite asserts the round trip, including correlation-ID
propagation from inbound metadata into the error payload.

\begin{minted}{python}
"""Server interceptor: map domain errors to gRPC status + rich details.

Rich details ride in the binary trailing-metadata key "error-details-bin",
the same transport trick Google APIs use for google.rpc.Status in
"grpc-status-details-bin". A client helper decodes them back.
"""
from __future__ import annotations

from typing import Any, Callable

import grpc

import telemetry_pb2
from telemetry_server import DomainError

TRAILER_KEY = "error-details-bin"


def _correlation_id(context: grpc.ServicerContext) -> str:
    for key, value in context.invocation_metadata():
        if key == "x-correlation-id":
            return str(value)
    return ""


def _abort_with_details(
    context: grpc.ServicerContext, err: DomainError
) -> None:
    details = telemetry_pb2.ErrorDetails(
        reason=err.reason,
        domain="telemetry.northwind.example",
        correlation_id=_correlation_id(context),
    )
    details.violations.extend(err.violations)
    context.set_trailing_metadata(
        ((TRAILER_KEY, details.SerializeToString()),)
    )
    context.abort(err.status_code, err.message)


def _wrap_unary(behavior: Callable) -> Callable:
    def handler(request: Any, context: grpc.ServicerContext) -> Any:
        try:
            return behavior(request, context)
        except DomainError as err:
            _abort_with_details(context, err)

    return handler


def _wrap_stream(behavior: Callable) -> Callable:
    def handler(request_or_iter: Any, context: grpc.ServicerContext) -> Any:
        try:
            yield from behavior(request_or_iter, context)
        except DomainError as err:
            _abort_with_details(context, err)

    return handler


class ErrorMappingInterceptor(grpc.ServerInterceptor):
    """Rebuilds each handler so DomainError becomes status + rich details."""

    def intercept_service(self, continuation, handler_call_details):
        handler = continuation(handler_call_details)
        if handler is None:
            return None
        deserializer = handler.request_deserializer
        serializer = handler.response_serializer
        if handler.request_streaming and handler.response_streaming:
            return grpc.stream_stream_rpc_method_handler(
                _wrap_stream(handler.stream_stream),
                request_deserializer=deserializer,
                response_serializer=serializer,
            )
        if handler.request_streaming:
            return grpc.stream_unary_rpc_method_handler(
                _wrap_unary(handler.stream_unary),
                request_deserializer=deserializer,
                response_serializer=serializer,
            )
        if handler.response_streaming:
            return grpc.unary_stream_rpc_method_handler(
                _wrap_stream(handler.unary_stream),
                request_deserializer=deserializer,
                response_serializer=serializer,
            )
        return grpc.unary_unary_rpc_method_handler(
            _wrap_unary(handler.unary_unary),
            request_deserializer=deserializer,
            response_serializer=serializer,
        )


def decode_error_details(
    call_or_error: Any,
) -> telemetry_pb2.ErrorDetails | None:
    """Client side: pull ErrorDetails out of a failed call's trailers."""
    metadata = call_or_error.trailing_metadata() or ()
    for key, value in metadata:
        if key == TRAILER_KEY:
            details = telemetry_pb2.ErrorDetails()
            details.ParseFromString(value)
            return details
    return None
\end{minted}

\subsection{Listing 4 --- Client: Metadata, Deadline Budgets, Interceptor, Retry Policy}
\label{subsec:ch13-code-client}

The interceptor attaches \texttt{authorization} and \texttt{x-correlation-id}
to \emph{every} call, of every cardinality; the channel carries a retry policy
as a JSON \emph{service config} (bounded attempts, exponential backoff,
\texttt{UNAVAILABLE} only --- retrying a non-idempotent method that may have
reached the server is how you double-charge a customer).

\begin{minted}{python}
"""Client for the Northwind Robotics telemetry service (Chapter 13).

Shows metadata, per-call deadline budgets, an auth + correlation-ID
interceptor, and a retry policy installed via gRPC service config.
Offline-first: localhost only.
"""
from __future__ import annotations

import itertools
import json
from dataclasses import dataclass
from typing import Any, Callable, Iterable, Iterator

import grpc

import telemetry_pb2
import telemetry_pb2_grpc

# Retry policy delivered through the channel's service config. Retries are
# bounded (maxAttempts) and restricted to UNAVAILABLE -- the only status here
# that is provably safe to retry for these non-idempotent methods when the
# attempt never reached the server.
RETRY_SERVICE_CONFIG = json.dumps(
    {
        "methodConfig": [
            {
                "name": [{"service": "northwind.telemetry.v1.TelemetryService"}],
                "retryPolicy": {
                    "maxAttempts": 3,
                    "initialBackoff": "0.050s",
                    "maxBackoff": "0.500s",
                    "backoffMultiplier": 2.0,
                    "retryableStatusCodes": ["UNAVAILABLE"],
                },
                # Overall per-RPC timeout when the caller sets none.
                "timeout": "2.000s",
            }
        ]
    }
)


@dataclass(frozen=True)
class _ClientCallDetails:
    """grpc.ClientCallDetails is an interface; this is our concrete copy."""

    method: str
    timeout: float | None
    metadata: tuple[tuple[str, Any], ...]
    credentials: Any
    wait_for_ready: bool | None
    compression: Any


class AuthMetadataInterceptor(
    grpc.UnaryUnaryClientInterceptor,
    grpc.UnaryStreamClientInterceptor,
    grpc.StreamUnaryClientInterceptor,
    grpc.StreamStreamClientInterceptor,
):
    """Attaches auth and correlation metadata to every outgoing RPC."""

    def __init__(
        self, token: str, correlation_factory: Callable[[], str]
    ) -> None:
        self._token = token
        self._correlation_factory = correlation_factory

    def _augment(
        self, details: grpc.ClientCallDetails
    ) -> _ClientCallDetails:
        metadata = tuple(details.metadata or ()) + (
            ("authorization", f"Bearer {self._token}"),
            ("x-correlation-id", self._correlation_factory()),
        )
        return _ClientCallDetails(
            method=details.method,
            timeout=details.timeout,
            metadata=metadata,
            credentials=details.credentials,
            wait_for_ready=details.wait_for_ready,
            compression=details.compression,
        )

    def intercept_unary_unary(self, continuation, client_call_details, request):
        return continuation(self._augment(client_call_details), request)

    def intercept_unary_stream(self, continuation, client_call_details, request):
        return continuation(self._augment(client_call_details), request)

    def intercept_stream_unary(
        self, continuation, client_call_details, request_iterator
    ):
        return continuation(self._augment(client_call_details), request_iterator)

    def intercept_stream_stream(
        self, continuation, client_call_details, request_iterator
    ):
        return continuation(self._augment(client_call_details), request_iterator)


def make_channel(
    address: str,
    token: str = "synthetic-dev-token",
    correlation_factory: Callable[[], str] | None = None,
) -> grpc.Channel:
    """Channel with retries (service config) and the auth interceptor."""
    if correlation_factory is None:
        counter = itertools.count(1)

        def _next_id() -> str:
            return f"corr-{next(counter):06d}"

        correlation_factory = _next_id
    options = [
        ("grpc.enable_retries", 1),
        ("grpc.service_config", RETRY_SERVICE_CONFIG),
        ("grpc.keepalive_time_ms", 60_000),
    ]
    raw = grpc.insecure_channel(address, options=options)
    return grpc.intercept_channel(
        raw, AuthMetadataInterceptor(token, correlation_factory)
    )


def fetch_status(
    stub: telemetry_pb2_grpc.TelemetryServiceStub,
    robot_id: str,
    deadline_s: float = 1.0,
) -> telemetry_pb2.RobotStatusReply:
    """Unary call with an explicit deadline budget."""
    request = telemetry_pb2.RobotStatusRequest(robot_id=robot_id)
    return stub.GetRobotStatus(request, timeout=deadline_s)


def stream_frames(
    stub: telemetry_pb2_grpc.TelemetryServiceStub,
    robot_id: str,
    max_frames: int,
    deadline_s: float = 5.0,
) -> Iterator[telemetry_pb2.TelemetryFrame]:
    request = telemetry_pb2.StreamTelemetryRequest(
        robot_id=robot_id, max_frames=max_frames, period_ms=20
    )
    yield from stub.StreamTelemetry(request, timeout=deadline_s)


def main() -> None:
    import argparse

    parser = argparse.ArgumentParser(description="Telemetry gRPC client demo")
    parser.add_argument("--port", type=int, default=50051)
    args = parser.parse_args()
    channel = make_channel(f"localhost:{args.port}")
    stub = telemetry_pb2_grpc.TelemetryServiceStub(channel)

    reply = fetch_status(stub, "RB-1001", deadline_s=1.0)
    mode = telemetry_pb2.RobotMode.Name(reply.mode)
    print(f"status: {reply.robot_id} mode={mode} battery={reply.battery_pct:.2f}%")

    print("stream:")
    for frame in stream_frames(stub, "RB-1001", max_frames=3):
        print(f"  frame battery={frame.battery_pct:.2f} mode={frame.mode}")

    frames = [
        telemetry_pb2.TelemetryFrame(robot_id="RB-1001", battery_pct=b)
        for b in (91.5, 90.8, 88.2)
    ]
    summary = stub.UploadTelemetry(iter(frames), timeout=2.0)
    print(f"upload: accepted={summary.accepted} rejected={summary.rejected}")
    channel.close()


if __name__ == "__main__":
    main()
\end{minted}

\subsection{Listing 5 --- Bidirectional Streaming with Flow-Control-Aware Production}
\label{subsec:ch13-code-bidi}

A producer thread writes commands into a \emph{bounded} queue; the gRPC request
generator drains it. When the window is full, \texttt{put()} blocks ---
application-level backpressure stacked on top of HTTP/2's per-stream flow
control. This is the canonical shape for any bidi producer fed from a source
that outruns the network.

\begin{minted}{python}
"""Bidirectional control-loop demo with flow-control-aware production.

A producer thread writes commands into a bounded queue; the gRPC request
generator drains it. When the queue is full the producer blocks, giving the
application its own backpressure on top of HTTP/2 flow control. Offline:
localhost only, synthetic commands, deterministic pacing.
"""
from __future__ import annotations

import queue
import threading
import time
from typing import Iterator

import grpc

import telemetry_pb2
import telemetry_pb2_grpc

_DONE: object = object()  # sentinel closing the command stream


def _producer(
    outbox: queue.Queue,
    robot_id: str,
    count: int,
    interval_s: float,
) -> None:
    try:
        for seq in range(count):
            if seq % 5 == 4:
                command = telemetry_pb2.ControlCommand(
                    robot_id=robot_id, sequence=seq, estop=False
                )
            else:
                command = telemetry_pb2.ControlCommand(
                    robot_id=robot_id,
                    sequence=seq,
                    set_mode=telemetry_pb2.ROBOT_MODE_MOVING,
                )
            # put() blocks while the window is full: application backpressure.
            outbox.put(command)
            time.sleep(interval_s)
    finally:
        outbox.put(_DONE)


def _commands(
    inbox: queue.Queue,
) -> Iterator[telemetry_pb2.ControlCommand]:
    while True:
        item = inbox.get()
        if item is _DONE:
            return
        yield item


def run_control_loop(
    stub: telemetry_pb2_grpc.TelemetryServiceStub,
    robot_id: str,
    count: int,
    interval_s: float = 0.005,
    window: int = 8,
    deadline_s: float = 10.0,
) -> list[telemetry_pb2.TelemetryFrame]:
    """Drive a bidi ControlLoop; returns every frame the server echoed."""
    outbox: queue.Queue = queue.Queue(maxsize=window)
    producer = threading.Thread(
        target=_producer,
        args=(outbox, robot_id, count, interval_s),
        daemon=True,
    )
    producer.start()
    frames: list[telemetry_pb2.TelemetryFrame] = []
    try:
        for frame in stub.ControlLoop(_commands(outbox), timeout=deadline_s):
            frames.append(frame)
    finally:
        producer.join(timeout=2.0)
    return frames


def main() -> None:
    import argparse

    from telemetry_client import make_channel

    parser = argparse.ArgumentParser(description="Bidirectional control-loop demo")
    parser.add_argument("--port", type=int, default=50051)
    args = parser.parse_args()
    channel = make_channel(f"localhost:{args.port}")
    stub = telemetry_pb2_grpc.TelemetryServiceStub(channel)
    frames = run_control_loop(stub, robot_id="RB-1001", count=10)
    for frame in frames:
        mode = telemetry_pb2.RobotMode.Name(frame.mode)
        print(f"ack robot={frame.robot_id} mode={mode} t={frame.observed_at.seconds}")
    channel.close()


if __name__ == "__main__":
    main()
\end{minted}

\subsection{Listing 6 --- Offline Test Suite: Deadline, Cancellation, and Stream Error Paths}
\label{subsec:ch13-code-tests}

The fixture binds the in-process server to port \texttt{0} (the OS assigns a
free \texttt{localhost} port) --- no fixed port collisions, no network, no
credentials. Eight tests cover: unary happy path, rich-error round trip,
deadline exceeded against the deliberately slow robot, server-stream bounds,
client cancellation observed as \texttt{CANCELLED}, mid-stream validation
abort, and the bidi round trip with fully deterministic timestamp assertions.

\begin{minted}{python}
"""Offline pytest suite for the Chapter 13 telemetry service.

The server runs in-process on an ephemeral localhost port (port 0); no
network, no credentials, deterministic data. Covers all four RPC modes,
deadline propagation, cancellation, and stream error paths.
"""
from __future__ import annotations

import grpc
import pytest

import telemetry_pb2
import telemetry_pb2_grpc
from bidi_demo import run_control_loop
from error_middleware import ErrorMappingInterceptor, decode_error_details
from telemetry_client import make_channel
from telemetry_server import BASE_EPOCH_S, make_server


@pytest.fixture()
def stub():
    server = make_server(interceptors=(ErrorMappingInterceptor(),))
    port = server.add_insecure_port("localhost:0")
    server.start()
    channel = make_channel(
        f"localhost:{port}", correlation_factory=lambda: "corr-test-0001"
    )
    yield telemetry_pb2_grpc.TelemetryServiceStub(channel)
    channel.close()
    server.stop(None)


def _frame(robot_id: str, battery: float) -> telemetry_pb2.TelemetryFrame:
    frame = telemetry_pb2.TelemetryFrame(robot_id=robot_id, battery_pct=battery)
    frame.observed_at.FromSeconds(BASE_EPOCH_S)
    return frame


def test_unary_status_happy_path(stub) -> None:
    reply = stub.GetRobotStatus(
        telemetry_pb2.RobotStatusRequest(robot_id="RB-1001"), timeout=5.0
    )
    assert reply.robot_id == "RB-1001"
    assert reply.mode == telemetry_pb2.ROBOT_MODE_MOVING
    assert 0.0 <= reply.battery_pct <= 100.0
    assert reply.last_frame.joint_positions_rad["shoulder"] != 0.0


def test_unknown_robot_maps_to_not_found_with_rich_details(stub) -> None:
    with pytest.raises(grpc.RpcError) as excinfo:
        stub.GetRobotStatus(
            telemetry_pb2.RobotStatusRequest(robot_id="RB-9999"), timeout=5.0
        )
    err = excinfo.value
    assert err.code() == grpc.StatusCode.NOT_FOUND
    details = decode_error_details(err)
    assert details is not None
    assert details.reason == "ROBOT_NOT_FOUND"
    assert details.correlation_id == "corr-test-0001"


def test_deadline_exceeded_propagates(stub) -> None:
    # RB-SLOW stalls 200 ms server-side; the 50 ms client deadline wins.
    with pytest.raises(grpc.RpcError) as excinfo:
        stub.GetRobotStatus(
            telemetry_pb2.RobotStatusRequest(robot_id="RB-SLOW"), timeout=0.05
        )
    assert excinfo.value.code() == grpc.StatusCode.DEADLINE_EXCEEDED


def test_server_stream_respects_max_frames(stub) -> None:
    request = telemetry_pb2.StreamTelemetryRequest(
        robot_id="RB-1002", max_frames=4, period_ms=5
    )
    frames = list(stub.StreamTelemetry(request, timeout=5.0))
    assert len(frames) == 4
    assert [f.observed_at.seconds for f in frames] == [
        BASE_EPOCH_S + i * 5 for i in range(4)
    ]


def test_client_cancellation_stops_server_stream(stub) -> None:
    request = telemetry_pb2.StreamTelemetryRequest(
        robot_id="RB-1001", max_frames=100_000, period_ms=5
    )
    call = stub.StreamTelemetry(request, timeout=30.0)
    first = next(call)
    assert first.robot_id == "RB-1001"
    call.cancel()
    with pytest.raises(grpc.RpcError) as excinfo:
        next(call)
    assert excinfo.value.code() == grpc.StatusCode.CANCELLED


def test_upload_rejects_out_of_range_battery(stub) -> None:
    frames = [_frame("RB-1001", 90.0), _frame("RB-1001", 140.0), _frame("RB-1001", 50.0)]
    call = stub.UploadTelemetry.future(iter(frames), timeout=5.0)
    summary = call.result()
    assert (summary.accepted, summary.rejected) == (2, 1)
    assert summary.min_battery_pct == pytest.approx(50.0)
    assert summary.max_battery_pct == pytest.approx(90.0)


def test_upload_unknown_robot_aborts_mid_stream_with_details(stub) -> None:
    frames = [_frame("RB-1001", 90.0), _frame("RB-9999", 10.0)]
    call = stub.UploadTelemetry.future(iter(frames), timeout=5.0)
    with pytest.raises(grpc.RpcError) as excinfo:
        call.result()
    err = excinfo.value
    assert err.code() == grpc.StatusCode.INVALID_ARGUMENT
    details = decode_error_details(err)
    assert details is not None
    assert details.reason == "FRAME_VALIDATION_FAILED"
    assert details.violations[0].field == "robot_id"
    assert details.correlation_id == "corr-test-0001"


def test_control_loop_bidi_round_trip(stub) -> None:
    frames = run_control_loop(
        stub, robot_id="RB-1001", count=10, interval_s=0.001, deadline_s=10.0
    )
    assert len(frames) == 10
    # Sequence numbers echo through to deterministic timestamps.
    assert [f.observed_at.seconds for f in frames] == [
        BASE_EPOCH_S + i * 5 for i in range(10)
    ]
    # seq % 5 == 4 carries estop=False -> IDLE; everything else MOVING.
    assert frames[4].mode == telemetry_pb2.ROBOT_MODE_IDLE
    assert frames[0].mode == telemetry_pb2.ROBOT_MODE_MOVING
\end{minted}

Verified output (Python 3.13, grpcio 1.83): all eight tests pass, and the
client and bidi demos print the deterministic battery and mode sequences shown
in the lab (Section~\ref{sec:ch13-lab}).

%% ----------------------------------------------------------------------------
\section{Common Pitfalls \& Anti-Patterns}
\label{sec:ch13-pitfalls}
\begin{enumerate}
  \item \textbf{Missing deadlines.} The default gRPC deadline is
        \emph{infinity}. A client without \texttt{timeout=} will wait forever
        on a wedged server, and the leaked threads/sockets accumulate into an
        outage. Rule: \emph{every call site sets a deadline}, and servers
        re-check \texttt{time\_remaining()} after any blocking step. gRPC's
        infinity default is the single most dangerous default in the stack.
  \item \textbf{Ignoring cancellation.} A server-streaming handler that never
        polls \texttt{context.is\_active()} keeps serializing and sending to a
        client that left minutes ago, burning CPU and holding locks. Treat
        \texttt{is\_active()} checks as loop hygiene, exactly like the
        \texttt{StreamTelemetry} listing.
  \item \textbf{The L4 load-balancer hotspot.} Putting a classic TCP load
        balancer in front of gRPC and declaring victory. Connections, not
        RPCs, get balanced; see the war story in Section~\ref{sec:ch13-lab}.
        Mitigate with client-side \texttt{round\_robin}, xDS/lookaside, an L7
        proxy, or at minimum \texttt{max\_connection\_age} so clients redial.
  \item \textbf{Unbounded streams.} \texttt{rpc Watch(X) returns (stream Y)}
        with no pagination, no max count, no heartbeat, and a consumer slower
        than the producer. HTTP/2 flow control will pause the sender, and your
        server's send buffers grow until OOM. Bound every stream
        (\texttt{max\_frames} in our IDL), add application-level backpressure
        (the bounded queue in Listing~5), and set deadlines on streams too.
  \item \textbf{Leaking proto field numbers.} Deleting a field and letting the
        next engineer reuse its tag. A stale peer then decodes old bytes as the
        new field --- silent corruption, no error. Always \texttt{reserve} the
        tag and name first (Listing~1), and review \texttt{.proto} diffs like
        schema migrations (Chapter~4).
  \item \textbf{Catching \texttt{grpc.RpcError} without reading
        \texttt{.code()}.} \texttt{except grpc.RpcError: retry()} conflates
        INVALID\_ARGUMENT (your bug --- retrying is a hot loop) with
        UNAVAILABLE (transient --- retrying is correct). Branch on
        \texttt{err.code()}; consult Table~\ref{tab:ch13-status}; only
        UNAVAILABLE (and ABORTED for transactional ops) is generically
        retryable.
  \item \textbf{Treating gRPC as a drop-in REST replacement for public APIs.}
        No browser support without a proxy, no human-readable payloads, no CDN
        caching, unfamiliar error surface for third parties, and payload
        inspection tooling that assumes JSON. Keep public edges on HTTP/JSON
        (transcode if you must share a contract); use gRPC where you control
        both ends.
  \item \textbf{Per-call channels.} Creating a channel (TCP + HTTP/2 setup)
        per RPC throws away multiplexing and adds handshake latency to every
        call. Channels are heavyweight and thread-safe: build one per target,
        share it, close it at shutdown.
  \item \textbf{Retried streaming calls without thinking.} A retry policy on a
        streaming RPC replays the request iterator --- fine for idempotent
        reads, catastrophic for ``append these events.'' Scope
        \texttt{retryPolicy} per method in the service config, and prefer
        hedging/retries only on methods you have marked idempotent.
  \item \textbf{Overstuffed metadata.} Metadata is HPACK-compressed but not
        free; multi-kilabyte headers on every call defeat the point. Keep
        metadata to credentials, correlation IDs, and small routing hints ---
        payloads belong in messages.
\end{enumerate}

%% ----------------------------------------------------------------------------
\section{Hands-On Production Lab \& Real-World Case Study}
\label{sec:ch13-lab}

\subsection{Lab: The Telemetry Service End to End}

All commands are PowerShell, offline, from a working directory containing the
six listings.

\textbf{Step 1 --- environment and codegen.}

\begin{minted}{text}
python -m venv .venv
.\.venv\Scripts\Activate.ps1
pip install grpcio grpcio-tools protobuf pytest
python -m grpc_tools.protoc -I. --python_out=. --grpc_python_out=. telemetry.proto
\end{minted}

\textbf{Step 2 --- run the suite.} \texttt{python -m pytest -q} should report
\texttt{8 passed}. Read the tests as executable specifications of the deadline,
cancellation, and error-path semantics.

\textbf{Step 3 --- exercise all four modes against a live server.} In one
terminal: \texttt{python telemetry\_server.py --port 50051}. In a second
terminal: \texttt{python telemetry\_client.py --port 50051} (unary +
server-streaming + client-streaming) and \texttt{python bidi\_demo.py --port
50051} (bidirectional). Expected client output:

\begin{minted}{text}
status: RB-1001 mode=ROBOT_MODE_MOVING battery=95.00%
stream:
  frame battery=95.00 mode=2
  frame battery=94.97 mode=2
  frame battery=94.96 mode=2
upload: accepted=3 rejected=0
\end{minted}

\textbf{Step 4 --- grpcurl (optional but instructive).} Reflection requires
\texttt{pip install grpcio-reflection} and two lines of server setup
(\texttt{reflection.enable\_server\_reflection}); then:

\begin{minted}{text}
grpcurl -plaintext localhost:50051 list
grpcurl -plaintext -d "{\"robot_id\": \"RB-1001\"}" localhost:50051 `
  northwind.telemetry.v1.TelemetryService/GetRobotStatus
grpcurl -plaintext -d "{\"robot_id\": \"RB-9999\"}" localhost:50051 `
  northwind.telemetry.v1.TelemetryService/GetRobotStatus   # NOT_FOUND (code 5)
\end{minted}

\textbf{Step 5 --- observe deadline propagation.} Call the slow robot with a
generous then a tight budget and compare outcomes:

\begin{minted}{python}
import grpc, telemetry_pb2, telemetry_pb2_grpc
channel = grpc.insecure_channel("localhost:50051")
stub = telemetry_pb2_grpc.TelemetryServiceStub(channel)
req = telemetry_pb2.RobotStatusRequest(robot_id="RB-SLOW")
print(stub.GetRobotStatus(req, timeout=5.0).robot_id)   # succeeds: budget ample
try:
    stub.GetRobotStatus(req, timeout=0.05)               # 50 ms < 200 ms stall
except grpc.RpcError as err:
    print(err.code())                                    # StatusCode.DEADLINE_EXCEEDED
\end{minted}

The server-side \texttt{\_require\_budget} re-check after the stall is what
turns ``slow'' into a clean \texttt{DEADLINE\_EXCEEDED} instead of a late,
wasted response.

\subsection{War Story: The L4 Hotspot That Pinned All Traffic to One Pod}

A logistics company (thinly fictionalized) ran a gRPC telemetry ingestion
fleet behind a cloud Layer-4 load balancer: 20 pods, autoscaled, each client
gateway holding one connection per target. Alerting fired at 03:40: p99
latency climbing on \emph{one} pod, CPU at 98\%, while the other nineteen
idled at 8\%. The autoscaler, watching \emph{fleet-average} CPU at 12\%,
politely did nothing.

The mechanism was exactly Section~\ref{subsec:ch13-operational}'s L4 problem,
with a queue-theoretic twist. The ingest clients were long-lived sidecars that
dialed once at boot and kept the connection for days. All twelve sidecars had
booted during a deploy when pod \texttt{telemetry-7} was the only ready
endpoint behind the LB --- so all twelve connections, and therefore
\emph{every RPC}, terminated on that one pod. Kubernetes service load
balancing had done its job perfectly: it balanced \emph{connections} (twelve
of them, all during a one-pod window) and had no idea that each connection
carried thousands of RPCs per second.

The on-call runbook failed twice. Restarting \texttt{telemetry-7} moved the
hotspot: clients re-dialed, the LB again picked the lowest-connection backend
--- which was the freshly restarted pod --- and within ninety seconds the new
pod was the hotspot. Scaling from 20 to 30 pods did nothing at all, because no
client ever re-dialed to discover the new backends. The incident lasted four
hours and its ``fix'' --- restarting all the sidecars simultaneously to force
mass redial --- was indistinguishable from an attack on their own fleet.

The durable fixes, in order of impact: \textbf{(1)} server-side
\texttt{max\_connection\_age} of five minutes with a thirty-second grace, so
every connection is periodically and gently recycled across the fleet ---
this alone would have capped the incident at minutes; \textbf{(2)} client-side
\texttt{round\_robin} with DNS resolution of all pod IPs, moving balance from
per-connection to per-RPC; \textbf{(3)} alerting on \emph{per-pod} p99 and
RPC-count skew (a $10\times$ max/median ratio), not fleet averages. Postmortem
quote worth memorizing: \emph{``We load-balanced the connections and called it
load balancing. Nobody owned the RPCs.''}

%% ----------------------------------------------------------------------------
\section{Self-Assessment Exercises \& Deep-Dive Questions}
\label{sec:ch13-exercises}

\begin{enumerate}
  \item \textbf{Wire decoding.} A server receives the bytes
        \texttt{00 00 00 00 0A} followed by ten protobuf bytes on stream 3.
        State the compression flag, message length, and what the receiver does
        if the stream ends after six of the ten payload bytes.
  \item \textbf{Framing arithmetic.} A response message serializes to 300
        bytes. Write out the exact 5-byte prefix in hex. What is the largest
        gRPC message this framing can express, and what default limit does
        grpc-python actually enforce (look up \texttt{grpc.max\_receive\_message\_length})?
  \item \textbf{Mode selection.} For each, pick a mode and defend it:
        (a)~pushing firmware-update progress to an operator console;
        (b)~uploading a day of buffered sensor logs; (c)~a trader-style
        command/ack console; (d)~fetching one robot's current pose.
  \item \textbf{Deadline math.} A client sets a 500 ms deadline. The proxy hop
        consumes 40 ms. What \texttt{grpc-timeout} value does the downstream
        server see (approximately), and in what unit suffix could it be
        encoded? What happens if the proxy \emph{adds} 50 ms of budget instead?
  \item \textbf{Retry safety.} The retry policy in Listing~4 names only
        \texttt{UNAVAILABLE}. Explain precisely why retrying
        \texttt{DEADLINE\_EXCEEDED} is dangerous for \texttt{UploadTelemetry}
        but harmless for \texttt{GetRobotStatus}.
  \item \textbf{Evolution drill.} Product wants to replace
        \texttt{battery\_pct} (tag 3, \texttt{double}) with
        \texttt{battery\_millivolts} (\texttt{uint32}). Write the safe two-phase
        migration using \texttt{reserved} and dual-publish, referencing
        Chapter~4's compatibility classes.
  \item \textbf{Interceptor design.} Sketch a server interceptor that enforces
        a per-token rate limit (Chapter~14 previews the algorithms) and returns
        \texttt{RESOURCE\_EXHAUSTED} with a \texttt{RetryInfo}-style detail in
        a \texttt{-bin} trailer. Which of the four handler shapes need
        different treatment and why?
  \item \textbf{Debug the symptom.} Clients see \texttt{UNAVAILABLE} spikes
        every five minutes, synchronized with nothing in the deploy calendar.
        Given the server options in Listing~2, form a hypothesis and the
        one-line config change that confirms it.
  \item \textbf{Transcoding design.} You must expose \texttt{StreamTelemetry}
        to a browser dashboard. Compare grpc-web vs.\ JSON transcoding vs.\ a
        bespoke WebSocket shim on contract fidelity, ops cost, and browser
        support; pick one and justify.
  \item \textbf{Deep dive.} Explain why \texttt{grpc-status} lives in trailers
        even for unary calls, what a trailers-only response is, and which
        intermediary behaviors (caches, L4 proxies, browsers) each property
        enables or forbids.
\end{enumerate}
%% ----------------------------------------------------------------------------
\section{Interview Q\&A Vault}
\label{sec:ch13-qa}

Ordered junior $\rightarrow$ staff. Answers are calibrated to what a strong
candidate says in two to six sentences.

\begin{enumerate}[label=\textbf{Q\arabic*.}, leftmargin=1.5cm]
  % --- junior ---
  \item \textbf{What is an RPC, and what are the stub and skeleton?}
        A remote procedure call is a network invocation dressed as a local
        function call. The \emph{stub} is client-side generated code exposing
        the typed signature and marshalling arguments to bytes; the
        \emph{skeleton} (servicer) is server-side generated glue that
        unmarshals and calls your implementation. Both derive from one IDL
        file, so the two ends cannot disagree about types.
  \item \textbf{What two technologies does gRPC combine, and why?}
        Protocol Buffers as IDL and binary encoding (compact, typed,
        evolvable) and HTTP/2 as transport (multiplexed streams, HPACK header
        compression, flow control). Together they give small payloads, cheap
        concurrency, and generated contracts~\cite{grpcdocs,protobufdocs}.
  \item \textbf{Explain the 5-byte gRPC message prefix.}
        Every message on a stream is length-prefixed: byte 0 is a compression
        flag (0/1), bytes 1--4 are the big-endian uint32 message length, then
        the serialized protobuf payload. HTTP/2 DATA frame boundaries are
        irrelevant to message boundaries --- the prefix is the delimiter, which
        is what lets one stream carry arbitrarily many ordered messages.
  \item \textbf{What does the \texttt{:path} pseudo-header look like for a
        gRPC call?}
        \texttt{/}\emph{package.Service}\texttt{/}\emph{Method}, e.g.\
        \texttt{/northwind.telemetry.v1.TelemetryService/GetRobotStatus}. The
        method is always POST; routing keys entirely off \texttt{:path}.
  \item \textbf{What \texttt{content-type} does gRPC use, and what HTTP status
        does a failed gRPC call return?}
        \texttt{application/grpc} (often \texttt{application/grpc+proto}). The
        HTTP status is 200 even for failed calls --- the outcome lives in the
        \texttt{grpc-status} trailer, because for streaming calls the failure
        may be discovered after response bytes have flowed. Interviewers ask
        this to see if you have ever looked at the wire.
  \item \textbf{What is a trailers-only response?}
        When a call fails before producing any message, the server may send a
        single HEADERS frame containing \texttt{:status 200},
        \texttt{grpc-status}, and END\_STREAM --- headers and trailers
        collapsed into one frame. Most gRPC error responses look like this on
        the wire.
  \item \textbf{Name the four RPC modes and a use case for each.}
        Unary (request/reply --- \texttt{GetRobotStatus}); server streaming
        (subscriptions, feeds --- \texttt{StreamTelemetry}); client streaming
        (bulk ingestion --- \texttt{UploadTelemetry}); bidirectional streaming
        (interactive control sessions --- \texttt{ControlLoop}).
  \item \textbf{Why does proto3 have no \texttt{required} fields?}
        Because required fields can never be safely removed: every old reader
        rejects messages missing them, so the field is load-bearing forever.
        Proto3 makes everything optional-with-defaults and moves validation to
        application logic, which can evolve.
  \item \textbf{What is a \texttt{oneof}?}
        A set of fields of which at most one is set; setting one clears the
        others, and presence is preserved on the wire so the receiver can ask
        \texttt{WhichOneof}. Use it for genuine alternatives, and note the
        evolution trap: moving an existing field into a \texttt{oneof} changes
        its wire semantics.
  \item \textbf{Why must protobuf tag numbers never be reused?}
        The tag, not the name, is the field's identity on the wire. A stale
        peer holding old bytes will decode them into the new field silently ---
        wrong data, no error. Retire tags with \texttt{reserved} so the
        compiler enforces the ban~\cite{protobufdocs}.

  % --- mid ---
  \item \textbf{Deadline vs.\ timeout --- is there a difference in gRPC?}
        In gRPC they are one mechanism seen from two ends: the client sets a
        timeout (relative), the core converts it to an absolute deadline, and
        transmits the \emph{remaining} budget in \texttt{grpc-timeout}. The
        server reads it via \texttt{context.time\_remaining()}. Because only
        the remainder travels, a multi-hop call tree shares one end-to-end
        budget --- that is deadline propagation.
  \item \textbf{Walk through deadline propagation across three hops.}
        Client budgets 500 ms; hop 1 receives it, spends 40 ms, and calls hop 2
        with \texttt{timeout=context.time\_remaining()} $\approx$ 460 ms; hop 2
        repeats to hop 3. Any hop may also cap the budget tighter (a stricter
        local policy), but must never widen it. When the budget expires, the
        innermost straggler surfaces DEADLINE\_EXCEEDED and every enclosing
        call unwinds with the same code.
  \item \textbf{Per-attempt vs.\ overall deadline --- how do retries
        interact?}
        All retry attempts draw from the same overall deadline, spaced by the
        configured backoff. If the overall deadline is smaller than
        (attempts $\times$ expected attempt time + backoff), later attempts are
        dead on arrival. Set overall deadlines from the p99 of a \emph{full}
        retry sequence, and remember gRPC only retries when it can prove the
        RPC never reached the server logic (or the method is idempotent).
  \item \textbf{What is gRPC metadata and what is the \texttt{-bin} suffix?}
        Metadata is the name/value pairs carried in HTTP/2 header and trailer
        blocks --- the gRPC equivalent of HTTP headers. Keys are lowercase
        ASCII; keys ending in \texttt{-bin} carry arbitrary binary (base64 on
        the wire, raw bytes in the API). Rich error details and auth tokens
        ride here; \texttt{grpc-*} keys are reserved.
  \item \textbf{How do you attach authentication to gRPC calls?}
        As metadata --- conventionally \texttt{authorization: Bearer <token>}
        --- applied by a \emph{client interceptor} so no call site can forget
        it, over TLS in production. The server validates in a \emph{server
        interceptor} and rejects with UNAUTHENTICATED before business logic
        runs. Channel credentials (TLS) secure the connection; call credentials
        (the token) authenticate the RPC; gRPC composes the two.
  \item \textbf{What is an interceptor and where would you use one?}
        Middleware for gRPC: client interceptors wrap outbound calls (auth,
        correlation IDs, client metrics, hedging), server interceptors wrap
        inbound handlers (authz, logging, error mapping, rate limiting). They
        come in unary and stream flavors for each side; streaming interceptors
        wrap iterators rather than single messages. Listing~4 and Listing~3 are
        working examples of both sides.
  \item \textbf{How do you configure retries in gRPC?}
        Via a JSON \emph{service config} on the channel:
        \texttt{retryPolicy} with \texttt{maxAttempts}, exponential backoff
        (\texttt{initialBackoff}, \texttt{backoffMultiplier},
        \texttt{maxBackoff}), and \texttt{retryableStatusCodes} --- typically
        only UNAVAILABLE. Hedging (\texttt{hedgingPolicy}) is the latency
        twin: fire a duplicate after a delay, take whichever answers first.
        Both need server-side throttling awareness (\texttt{retryThrottling})
        to avoid retry storms.
  \item \textbf{Which status codes are safe to retry?}
        UNAVAILABLE is the canonical retryable code (transient, and the
        framework retries only when the attempt provably never ran). ABORTED is
        retryable for transactional conflicts at the application level.
        RESOURCE\_EXHAUSTED after the server-supplied delay. Everything else ---
        INVALID\_ARGUMENT, NOT\_FOUND, PERMISSION\_DENIED, DEADLINE\_EXCEEDED
        --- retrying unchanged is a bug amplifier.
  \item \textbf{INVALID\_ARGUMENT vs.\ FAILED\_PRECONDITION vs.\
        OUT\_OF\_RANGE?}
        INVALID\_ARGUMENT: the request is malformed regardless of system state
        (bad enum, missing field). FAILED\_PRECONDITION: the request is fine
        but system state forbids it (deleting a non-empty warehouse zone).
        OUT\_OF\_RANGE: a value lies outside a defined valid interval (seeking
        past the end of a stream). The distinctions matter because clients
        automate on them: only FAILED\_PRECONDITION hints ``fix state, then
        retry.''
  \item \textbf{UNAUTHENTICATED vs.\ PERMISSION\_DENIED?}
        UNAUTHENTICATED (16): credentials missing, expired, or invalid ---
        re-authentication may succeed. PERMISSION\_DENIED (7): identity known,
        access refused --- re-auth will not help. Mapping to HTTP: 401 vs.\ 403.
        Clients key their token-refresh logic off this distinction, so getting
        it backwards causes refresh loops or stranded users.

  % --- senior ---
  \item \textbf{Design error semantics for a streaming RPC.}
        Three eras: errors \emph{before} the first message can be
        trailers-only with full details; errors \emph{mid-stream} terminate
        the stream with \texttt{grpc-status} + trailers, so the client must
        treat all received messages as a valid prefix and the status as
        authoritative; errors \emph{per item} that should not kill the stream
        must be modeled \emph{in the message} (an \texttt{oneof
        result \{T ok; ItemError err;\}}), not as statuses. Publish which
        statuses the method can return, attach rich details
        (\texttt{ErrorInfo}/\texttt{BadRequest}-style) in \texttt{-bin}
        trailers, and state idempotency so clients know whether resume is safe.
  \item \textbf{How do rich error details work on the wire?}
        A \texttt{google.rpc.Status} message --- code, message, and
        \texttt{repeated Any details} holding typed payloads like
        \texttt{BadRequest}, \texttt{RetryInfo}, \texttt{QuotaFailure} --- is
        serialized into the binary trailer \texttt{grpc-status-details-bin}.
        Clients unpack the typed details and act programmatically (field
        errors to forms, retry delays to backoff). Our chapter implements the
        identical pattern with a first-party \texttt{ErrorDetails} message and
        an \texttt{error-details-bin} trailer.
  \item \textbf{Why does gRPC need client-side load balancing?}
        Because gRPC multiplexes many RPCs over one long-lived HTTP/2
        connection, and L4 balancers balance \emph{connections}, not RPCs: one
        connection pins all its RPCs to one backend forever. Client-side LB
        (resolve all backends via DNS, pick per-RPC with \texttt{round\_robin})
        moves balancing to RPC granularity with no extra network hop. The
        alternatives are a lookaside control plane (grpclb), xDS from a mesh
        control plane, or terminating HTTP/2 at an L7 proxy that rebalances
        per stream~\cite{newman2021microservices}.
  \item \textbf{What is lookaside load balancing, and what is xDS?}
        Lookaside: the client queries a dedicated balancer service for a
        curated backend list, then connects directly --- control traffic via
        the balancer, data traffic point-to-point (gRPC's historical grpclb).
        xDS is the CNCF-standardized discovery API family from Envoy ---
        LDS/CDS/RDS/EDS --- through which a control plane pushes listeners,
        clusters, routes, and endpoints to data-plane clients; gRPC's proxyless
        service-mesh mode consumes xDS natively. xDS is where the industry
        converged; grpclb is legacy context.
  \item \textbf{How do \texttt{max\_connection\_age} and keepalive improve
        operability?}
        \texttt{max\_connection\_age} forces periodic reconnects so
        long-lived clients rediscover rebalanced backends --- the cheap fix
        for L4 pinning --- with a grace period to drain in-flight RPCs.
        Keepalive PINGs detect half-dead peers and keep NAT/firewall state
        alive on idle connections; client ping intervals must respect the
        server's minimum or the server answers with GOAWAY (``too many
        pings'').
  \item \textbf{How does flow control work for gRPC streams?}
        Two layers. HTTP/2 gives per-stream and per-connection credit windows:
        receivers grant bytes with WINDOW\_UPDATE, so a slow consumer pauses
        the sender's socket writes. gRPC adds application-level flow via your
        iterator pacing. The operational point: flow control converts
        slow-consumer problems into \emph{buffered-bytes} problems --- bound
        queues (Listing~5), bound stream lengths, and watch memory, not just
        throughput.
  \item \textbf{What happens on cancellation, end to end?}
        The client sends RST\_STREAM for the call's stream; the transport
        tears down stream state; the server's \texttt{context.is\_active()}
        flips false; well-written handlers observe it between iterations and
        return promptly. Cancellation is cooperative --- nothing forcibly
        kills your handler --- and the client observes code CANCELLED. Deadline
        expiry is the same machinery with DEADLINE\_EXCEEDED.
  \item \textbf{How would you test gRPC services without a network?}
        Bind the server to \texttt{localhost:0} (ephemeral port) inside a
        pytest fixture and exercise the real stack in-process, as Listing~6
        does --- you get genuine serialization, framing, deadlines, and
        cancellation with zero infrastructure. For pure unit tests, call the
        servicer methods with a fake \texttt{ServicerContext}. Reserve mocks of
        generated stubs for testing \emph{callers}, never the service.
  \item \textbf{What is server reflection and what does grpcurl need it for?}
        Reflection is an optional service (\texttt{grpc.reflection\ldots}) that
        lets clients enumerate services and fetch file descriptors at runtime.
        \texttt{grpcurl} uses it to build requests without a local
        \texttt{.proto}: \texttt{grpcurl -plaintext host:port list}. Enable it
        in dev and staging; consider disabling or authenticating it in
        production because it advertises your full API surface.
  \item \textbf{How do you expose a gRPC service to a browser?}
        Browsers cannot speak raw gRPC (no trailer access, no fine-grained
        HTTP/2 control), so: \textbf{grpc-web} --- a framed protocol variant
        through a proxy filter (Envoy), with generated TS stubs, supporting
        unary and server streaming; or \textbf{HTTP/JSON transcoding} ---
        \texttt{google.api.http} annotations on the proto, a gateway (Envoy
        transcoder, gRPC-Gateway) mapping REST paths to RPCs and statuses per
        Table~\ref{tab:ch13-status}. Transcoding doubles as your public REST
        facade over an internal gRPC core.

  % --- staff ---
  \item \textbf{When should a public API \emph{not} use gRPC?}
        When any of these dominate: browser-first consumers (grpc-web adds a
        proxy and drops bidi); third-party developers who expect
        curl-able JSON and familiar HTTP semantics; a need for CDN/edge
        caching, which is built on HTTP semantics gRPC bypasses; firewalled
        enterprise networks that pass 443/HTTP but break long-lived HTTP/2;
        or heavy reliance on web tooling (API gateways, WAFs, scanners) that
        speaks JSON. The staff answer adds: you can have both --- gRPC
        internally, transcode at the edge --- and the contract cost of
        maintaining that mapping is the real decision variable.
  \item \textbf{Design the load-balancing architecture for a 500-service
        gRPC mesh.}
        Layer it: (1) xDS from the mesh control plane for discovery, policy,
        and subsetting --- clients should not DNS-resolve 500 services
        themselves; (2) \texttt{round\_robin} (or least-request where
        available) per-RPC within a locality, with locality weighting to keep
        traffic in-zone; (3) server-side \texttt{max\_connection\_age} as the
        safety net for any pinned path; (4) outlier detection to eject sick
        backends; (5) per-pod skew metrics as SLIs. State explicitly what you
        are balancing (RPCs), what each layer owns, and how you would detect
        imbalance before users do.
  \item \textbf{A teammate proposes gRPC streaming to replace Kafka for event
        delivery. Evaluate.}
        Different guarantees. gRPC streams are ordered, flow-controlled, and
        low-latency, but ephemeral: no durable log, no replay, no consumer
        groups, no retention --- a disconnected consumer simply misses events.
        Streams are right for \emph{control-plane} and \emph{live-tail}
        traffic; the log is right for \emph{facts} you must not lose and must
        reprocess (Chapter~11). A defensible hybrid: Kafka for the durable
        spine, gRPC server-streaming as the fan-out edge --- with sequence
        numbers and a resume offset modeled in your own messages.
  \item \textbf{How do you version a gRPC API across an organization?}
        Prefer evolution over versioning: proto3's rules (add-only fields,
        \texttt{reserved} tags) let most changes ship in place
        (Chapter~9's discipline applies with tags instead of JSON keys). When a
        change is genuinely breaking, version the \emph{package}
        (\texttt{northwind.telemetry.v2}) so both generations coexist on the
        wire, run dual servers during migration, and retire v1 with measured
        deprecation. Package-level versions in the \texttt{:path} make routing,
        metrics, and sunset enforcement trivial --- this is why the version
        belongs in the package name from day one.
  \item \textbf{Debug: latency p99 doubles after moving gRPC behind a cloud
        L4 load balancer; CPU is low everywhere. Hypotheses?}
        Ordered by prior: (1) connection pinning plus \emph{unequal
        connection lifetimes} --- a few hot tenants' RPCs concentrate; check
        per-pod RPC counts. (2) The LB's idle timeout silently reaping
        connections without keepalive --- clients stall on dead sockets before
        reconnect; check for synchronized reconnect storms. (3) DNS TTL vs.\
        client resolver caching, so scale-outs are invisible for minutes. (4)
        MTU/fragmentation on the new path inflating DATA frame latency. The
        instrumentation answer: per-connection RPC counts, channel state
        transitions, and \texttt{grpc-timeout} distributions server-side.
  \item \textbf{Compare gRPC, REST+JSON, and GraphQL for an internal
        service-to-service call graph.}
        REST+JSON: universal tooling and debuggability, weakest contract and
        heaviest payload; fine at low call rates. GraphQL: client-driven
        selection kills over/under-fetching at the \emph{aggregation} edge, but
        its query planning cost and N+1 hazards make it wrong for
        machine-to-machine hot paths (Chapter~12). gRPC: strongest contract,
        smallest payload, first-class streaming, best fit for
        service-to-service --- at the price of human-unfriendly wire format and
        LB complexity. Staff-level answers pick per \emph{edge}, and can state
        the organization's default and its exceptions.
  \item \textbf{What would you standardize in a ``gRPC production checklist''
        for your company?}
        Deadlines on every call with propagation through
        \texttt{time\_remaining()}; \texttt{is\_active()} checks in every
        stream loop; bounded streams and bounded queues; a documented status
        code policy with rich details in \texttt{-bin} trailers; interceptors
        for auth/correlation/metrics, never hand-threaded metadata; retry
        policy per method with idempotency annotation; client-side LB or xDS
        plus \texttt{max\_connection\_age}; keepalive tuned to the middlebox
        reality; health-check service for every server; reflection and grpcurl
        in non-prod; contract CI that regenerates stubs and rejects
        tag-reuse diffs. That list \emph{is} this chapter, operationalized.
\end{enumerate}

%% ----------------------------------------------------------------------------
\section{External Bibliography \& Further Reading}
\label{sec:ch13-biblio}

\begin{itemize}
  \item \textbf{gRPC documentation} --- \texttt{https://grpc.io/docs/}: the
        authoritative guides for concepts, deadlines, metadata, interceptors,
        service config, load balancing, and the health-checking and reflection
        protocols~\cite{grpcdocs}.
  \item \textbf{gRPC status codes} --- \texttt{https://grpc.io/docs/guides/error/}
        and \texttt{google/rpc/code.proto}: the canonical code list and the
        HTTP mapping conventions used in Table~\ref{tab:ch13-status}.
  \item \textbf{Protocol Buffers documentation} ---
        \texttt{https://protobuf.dev/}: the proto3 language guide, well-known
        types, the wire format (pair with Chapter~4), and the officially
        documented evolution rules~\cite{protobufdocs}.
  \item \textbf{RFC 9113, \emph{HTTP/2}} --- streams, frames, flow control, and
        HPACK context that gRPC builds on; Chapter~1's companion
        reference~\cite{rfc9113}.
  \item \textbf{Kasun Indrasiri \& Danesh Kuruppu, \emph{gRPC: Up and Running}}
        (O'Reilly, 2020) --- the standard book-length treatment: IDL design,
        all four modes in depth, production concerns, and the gRPC ecosystem.
  \item \textbf{Google AIP (API Improvement Proposals)} ---
        \texttt{https://aip.dev/}: AIP-121 (resource-oriented design), AIP-122
        (resource names), AIP-158 (pagination), AIP-193 (errors) --- the design
        system behind Google's gRPC APIs; steal it wholesale for internal
        consistency.
  \item \textbf{CNCF xDS / Envoy documentation} ---
        \texttt{https://www.envoyproxy.io/docs/envoy/latest/api-docs/xds\_protocol}
        --- the discovery-API family behind proxyless service mesh and modern
        gRPC load balancing.
  \item \textbf{grpc-web} --- \texttt{https://github.com/grpc/grpc-web}: the
        browser bridge protocol and proxy model.
  \item \textbf{Martin Kleppmann, \emph{Designing Data-Intensive Applications}}
        (O'Reilly, 2017) --- the RPC fallacies, schema evolution, and
        distributed-systems framing underpinning this
        chapter~\cite{kleppmann2017ddia}.
  \item \textbf{Sam Newman, \emph{Building Microservices}, 2nd ed.}
        (O'Reilly, 2021) --- where RPC fits (and does not fit) among
        inter-service communication styles~\cite{newman2021microservices}.
\end{itemize}
%% ----------------------------------------------------------------------------
\section{Capstone Projects}
\label{sec:ch13-capstones}

All capstones are offline-first, use the Northwind Robotics universe, and run
on Windows/PowerShell with only this chapter's dependencies.

\subsection{Capstone 1 [junior] --- Unary Status Service with Deadline Discipline}

\textbf{Objective.} Stand up a single-RPC gRPC service and prove you understand
deadlines from both ends of the wire.

\textbf{Inputs/Datasets.} Listing~1's \texttt{telemetry.proto} reduced to
\texttt{GetRobotStatus} only; the synthetic fleet table
(\texttt{RB-1001}\ldots\texttt{RB-1003}, \texttt{RB-SLOW}).

\textbf{Steps.}
\begin{enumerate}
  \item Trim the proto, regenerate stubs, implement server and client.
  \item Set explicit deadlines at every call site; add the
        \texttt{\_require\_budget} server check.
  \item Demonstrate \texttt{DEADLINE\_EXCEEDED} against \texttt{RB-SLOW} with a
        50 ms client budget, and success with a 5 s budget.
\end{enumerate}

\textbf{Acceptance Criteria.}
\begin{itemize}
  \item[$\square$] \texttt{pytest} suite passes with happy-path, unknown-robot,
        and deadline-exceeded cases.
  \item[$\square$] No call site lacks a \texttt{timeout=}; no handler lacks a
        budget check.
  \item[$\square$] \texttt{grpcurl} (or a printed descriptor dump) shows the
        service shape.
\end{itemize}

\textbf{Stretch Goals.} Add \texttt{grpc.keepalive\_time\_ms} and log channel
state transitions; measure call latency at the client interceptor.

\subsection{Capstone 2 [mid] --- Server-Streaming Telemetry Monitor}

\textbf{Objective.} Build a resilient streaming consumer that respects
cancellation, bounds, and flow control.

\textbf{Inputs/Datasets.} \texttt{StreamTelemetry} from Listing~1; synthetic
frame generator from Listing~2.

\textbf{Steps.}
\begin{enumerate}
  \item Implement the server stream with \texttt{max\_frames},
        \texttt{period\_ms}, \texttt{is\_active()} checks.
  \item Build a CLI monitor that prints frames and cancels cleanly on
        KeyboardInterrupt (\texttt{call.cancel()}), verifying the server stops
        producing.
  \item Add a client-side rendering budget: drop-and-count frames when the
        renderer falls behind instead of buffering unboundedly.
\end{enumerate}

\textbf{Acceptance Criteria.}
\begin{itemize}
  \item[$\square$] Cancellation test passes: after \texttt{cancel()}, the
        server handler exits within one period.
  \item[$\square$] Server never emits more than \texttt{max\_frames}; client
        drop counter is observable.
  \item[$\square$] Deadline on the stream is enforced and tested.
\end{itemize}

\textbf{Stretch Goals.} Add a heartbeat frame every $n$ periods; detect a
stalled stream client-side via per-frame timeout.

\subsection{Capstone 3 [mid] --- Bulk Ingestion with Rich Validation Errors}

\textbf{Objective.} Client-streaming ingestion where bad data produces
\emph{actionable} errors, not stack traces.

\textbf{Inputs/Datasets.} \texttt{UploadTelemetry}, \texttt{ErrorDetails}, and
\texttt{FieldViolation} from Listing~1; a synthetic file of 1{,}000 frames,
5\% deliberately invalid (unknown robots, out-of-range battery).

\textbf{Steps.}
\begin{enumerate}
  \item Implement upload validation and the error-mapping interceptor
        (Listing~3).
  \item Extend \texttt{ErrorDetails} with \texttt{first\_bad\_index} so clients
        can resume after the failure point.
  \item Write the client-side decoder that renders violations as a table.
\end{enumerate}

\textbf{Acceptance Criteria.}
\begin{itemize}
  \item[$\square$] Mid-stream abort yields INVALID\_ARGUMENT plus decoded
        per-field violations and the correlation ID end-to-end.
  \item[$\square$] A fix-and-resume run accepts exactly the valid 95\%.
  \item[$\square$] All errors are asserted via \texttt{.code()}, never via
        message-string matching.
\end{itemize}

\textbf{Stretch Goals.} Port the details payload to
\texttt{google.rpc.Status} + \texttt{grpc-status-details-bin} using
\texttt{googleapis-common-protos}, and show the wire equivalence.

\subsection{Capstone 4 [senior] --- Bidirectional Control Plane with a Full
Interceptor Stack}

\textbf{Objective.} A production-shaped bidi service: windowed flow control,
auth + logging + metrics interceptors on both sides, and a retry policy that
provably cannot duplicate control commands.

\textbf{Inputs/Datasets.} \texttt{ControlLoop} (Listing~1), Listings~3--5; a
scripted scenario of 500 deterministic commands including estops.

\textbf{Steps.}
\begin{enumerate}
  \item Add sequence-number idempotency server-side (duplicate
        \texttt{sequence} $\Rightarrow$ replay cached ack).
  \item Implement a metrics interceptor (counts, latency histogram buckets per
        method) and a server authz interceptor rejecting a synthetic bad token
        with UNAUTHENTICATED.
  \item Build a two-hop deadline-propagation harness: a proxy servicer that
        forwards \texttt{ControlLoop} to a backend using
        \texttt{time\_remaining()}, and show a 300 ms end-to-end budget
        expiring cleanly at both hops.
\end{enumerate}

\textbf{Acceptance Criteria.}
\begin{itemize}
  \item[$\square$] 500-command run: exactly-once \emph{effects} despite forced
        client retries; ack order preserved.
  \item[$\square$] Metrics interceptor emits per-cardinality counters without
        touching business code.
  \item[$\square$] Deadline test shows the downstream hop observing a
        \emph{smaller} remaining budget than the client set.
  \item[$\square$] Bad-token calls fail before any servicer code runs.
\end{itemize}

\textbf{Stretch Goals.} Add hedging to unary reads only, and measure the p99
delta against the retry baseline with a deterministic latency-injection hook.

\subsection{Capstone 5 [staff] --- Edge Gateway: Transcoding plus Load-Balancing
Postmortem Lab}

\textbf{Objective.} Reproduce the war story of Section~\ref{sec:ch13-lab} in
miniature, fix it three ways, and expose the service safely to JSON clients ---
then write the incident report.

\textbf{Inputs/Datasets.} The full telemetry service; three server processes on
loopback ports simulating ``pods''; a connection-pinning harness that dials
through a single fixed address; a synthetic load script (fixed seed, fixed
schedule).

\textbf{Steps.}
\begin{enumerate}
  \item Demonstrate the hotspot: all clients dialed during a one-pod window;
        show per-pod RPC counts $10\times$ skewed.
  \item Apply, and measure separately: \texttt{max\_connection\_age}
        recycling; client-side \texttt{round\_robin} over all three pod
        addresses; and a least-loaded selection policy you design.
  \item Add a JSON transcoding shim (a small HTTP server mapping
        \texttt{GET /v1/robots/\{id\}/status} to \texttt{GetRobotStatus} and
        statuses per Table~\ref{tab:ch13-status}) with rich details rendered
        into an RFC-9457-shaped body.
  \item Write the postmortem: timeline, mechanism, why restart-and-scale
        made it worse, and the three fixes with measured before/after skew.
\end{enumerate}

\textbf{Acceptance Criteria.}
\begin{itemize}
  \item[$\square$] Pinned-phase skew $\geq 10\times$; each fix reduces max/median
        RPC ratio below $1.5\times$ in the harness.
  \item[$\square$] JSON edge returns correct HTTP codes for NOT\_FOUND,
        DEADLINE\_EXCEEDED, and INVALID\_ARGUMENT cases, with details intact.
  \item[$\square$] Postmortem identifies the balance unit (connections vs.\
        RPCs) as the root cause and states the per-pod SLI that would have
        paged earlier.
  \item[$\square$] Everything runs offline with \texttt{pytest -q} green and a
        single PowerShell driver script.
\end{itemize}

\textbf{Stretch Goals.} Replace the hand-rolled \texttt{round\_robin} with a
DNS-name-resolver-based config; add the \texttt{grpc.health.v1} service and
gate the shim's routing on \texttt{SERVING}.

%% ----------------------------------------------------------------------------
\section{Python Dependencies Appendix}
\label{sec:ch13-deps}

\subsection*{\texttt{requirements.txt}}

\begin{minted}{text}
grpcio>=1.62          # gRPC runtime: channels, server, interceptors, statuses
grpcio-tools>=1.62    # bundled protoc + Python/gRPC plugins for codegen
protobuf>=4.25        # message runtime used by generated *_pb2 modules
pytest>=8             # test runner for the offline suite
\end{minted}

Install from PowerShell (offline-friendly once wheels are cached):

\begin{minted}{text}
python -m venv .venv
.\.venv\Scripts\Activate.ps1
pip install -r requirements.txt
\end{minted}

\subsection{\texttt{grpcio}}

\textbf{Description.} The gRPC runtime for Python: channels and connections,
the threaded server, all four interceptor interfaces, status codes, metadata,
deadlines, and the service-config/retry machinery.

\textbf{Why included.} It is the subject of the chapter; every listing depends
on it directly.

\textbf{Alternatives.} \emph{Twirp} (Twitch) --- a simpler HTTP/1.1+JSON/proto
RPC with a fraction of the features (no bidi streaming, no HTTP/2 multiplexing)
and a fraction of the operational complexity. \emph{Connect-RPC} (Buf) ---
gRPC-compatible but also speaks plain HTTP+JSON, easing browser and curl access;
a strong modern choice when one contract must serve both worlds. \emph{Apache
Thrift} --- the previous generation: own transport, own IDL, mature, but no
HTTP/2 semantics and a smaller mindshare today. \emph{Plain HTTP/2} (e.g.\
\texttt{httpx} + handwritten framing) --- you reimplement stubs, deadlines, and
statuses yourself; educational, not advisable.

\textbf{Installation.} \texttt{pip install grpcio} (wheels for CPython 3.11+
on Windows; no compiler needed).

\subsection{\texttt{grpcio-tools}}

\textbf{Description.} Bundles the \texttt{protoc} compiler, the Python and
gRPC-Python plugins, and the well-known \texttt{.proto} files, exposed as
\texttt{python -m grpc\_tools.protoc}.

\textbf{Why included.} Code generation is core to the gRPC workflow; bundling
keeps the pipeline hermetic and offline --- no separate protoc download.

\textbf{Alternatives.} A standalone \texttt{protoc} binary plus
\texttt{grpc\_python\_plugin}; or the Buf toolchain (\texttt{buf generate}),
which adds linting and breaking-change detection and is the better choice at
organization scale.

\textbf{Installation.} \texttt{pip install grpcio-tools}; keep its version in
lockstep with \texttt{grpcio}.

\subsection{\texttt{protobuf}}

\textbf{Description.} The Protocol Buffers runtime that generated
\texttt{\_pb2} modules import: message classes, serialization, well-known
types, reflection over descriptors.

\textbf{Why included.} Transitive but load-bearing: pin it explicitly because
generated code validates its runtime version, and Chapter~4's hand-rolled codec
comparisons use it as the reference implementation.

\textbf{Alternatives.} None within the protobuf ecosystem (the runtime is
mandatory for generated code); the \emph{format-level} alternatives ---
MessagePack, Avro, CBOR --- are analyzed in Chapter~4.

\textbf{Installation.} \texttt{pip install protobuf}; match the major version
\texttt{grpcio-tools} generated against.

\subsection{\texttt{pytest}}

\textbf{Description.} The test runner and fixture system used for the offline
suite (Listing~6).

\textbf{Why included.} Consistent with every prior chapter; fixtures make the
in-process server pattern (bind port 0, yield stub, stop server) three lines
long.

\textbf{Alternatives.} \texttt{unittest} (stdlib, zero dependency --- the
fixture becomes \texttt{setUp/tearDown}); \texttt{nose2} (legacy, not
recommended).

\textbf{Installation.} \texttt{pip install pytest}.

%% ----------------------------------------------------------------------------
\section{Chapter Summary \& Key Takeaways}
\label{sec:ch13-summary}

\begin{itemize}
  \item gRPC = Protocol Buffers IDL + protobuf binary encoding + HTTP/2
        transport + generated stubs. The \texttt{.proto} is the contract; both
        ends are artifacts of it.
  \item On the wire, one RPC is one HTTP/2 stream; every message carries a
        5-byte prefix (1-byte compression flag + 4-byte big-endian length);
        \texttt{:path} is \texttt{/package.Service/Method}; the HTTP status is
        200 and the real outcome is the \texttt{grpc-status} trailer ---
        collapsed into a trailers-only HEADERS frame for early failures.
  \item Four modes --- unary, server streaming, client streaming,
        bidirectional --- are declared in the IDL and differ in cancellation
        and flow-control obligations; bound every stream and poll
        \texttt{is\_active()}.
  \item Deadlines propagate: \texttt{grpc-timeout} carries the
        \emph{remaining} budget hop to hop. Set a deadline on every call;
        re-check \texttt{time\_remaining()} after stalls; size overall budgets
        for the whole retry sequence.
  \item The 17 status codes are the error model; rich details travel in
        \texttt{-bin} trailers (\texttt{google.rpc.Status} or your own typed
        payload). Branch on \texttt{.code()}; retry only UNAVAILABLE-class
        failures of idempotent-or-unreached calls.
  \item Interceptors --- client and server, unary and stream --- are where
        auth, correlation IDs, logging, metrics, retry, and error mapping
        live. Business logic never threads metadata by hand.
  \item Operationally: L4 balancers pin gRPC to single backends (balance
        RPCs, not connections --- client-side LB, xDS, or
        \texttt{max\_connection\_age}); keepalive detects dead peers;
        reflection + \texttt{grpcurl} restore curl-ability; grpc-web or JSON
        transcoding bridge to browsers; \texttt{grpc.health.v1} is the
        readiness contract.
  \item gRPC is for hot internal paths between services you control. Public,
        browser-first, cache-dependent APIs remain HTTP/JSON territory ---
        transcode at the edge if one contract must serve both.
\end{itemize}

%% ----------------------------------------------------------------------------
\section{Quick-Reference Card}
\label{sec:ch13-quickref}

\subsection*{RPC Mode Selection}
\begin{table}[h]
  \centering
  \small
  \begin{tabularx}{\textwidth}{l X X}
    \toprule
    \textbf{Mode} & \textbf{Pick when} & \textbf{Watch out for} \\
    \midrule
    Unary & Single request/reply; command or query. & Missing deadlines; retry only if idempotent. \\
    Server streaming & Feeds, tailing, large incremental results. & Abandoned streams; bound with max count; check \texttt{is\_active()}. \\
    Client streaming & Bulk ingestion, uploads, aggregates. & Mid-stream validation must abort with details; buffer bounds. \\
    Bidirectional & Control sessions, interactive protocols. & No implicit req/resp pairing; app-level backpressure queue. \\
    \bottomrule
  \end{tabularx}
\end{table}

\subsection*{Status Codes You Must Know Cold}
\begin{table}[h]
  \centering
  \small
  \begin{tabularx}{\textwidth}{l X l}
    \toprule
    \textbf{Code} & \textbf{Meaning} & \textbf{Retry?} \\
    \midrule
    OK (0) & Success. & --- \\
    CANCELLED (1) & Caller cancelled. & Only if safe to redo. \\
    INVALID\_ARGUMENT (3) & Malformed request. & No. \\
    DEADLINE\_EXCEEDED (4) & Budget expired. & No (budget already spent). \\
    NOT\_FOUND (5) & Resource absent. & No. \\
    PERMISSION\_DENIED (7) & Known caller, refused. & No. \\
    RESOURCE\_EXHAUSTED (8) & Quota/rate limit. & After server delay. \\
    FAILED\_PRECONDITION (9) & State forbids op. & After fixing state. \\
    ABORTED (10) & Concurrency conflict. & Yes, whole operation. \\
    UNIMPLEMENTED (12) & Method unknown here. & No. \\
    INTERNAL (13) & Server bug. & No. \\
    UNAVAILABLE (14) & Transient unavailability. & Yes (canonical). \\
    UNAUTHENTICATED (16) & Bad/missing credentials. & After re-auth. \\
    \bottomrule
  \end{tabularx}
  \caption{Abridged; the full 17-code table with HTTP mappings is
  Table~\ref{tab:ch13-status}.}
\end{table}

\subsection*{Deadline \& Cancellation Checklist}
\begin{itemize}
  \item[$\square$] Every call site sets \texttt{timeout=}; no infinite defaults.
  \item[$\square$] Every hop forwards \texttt{context.time\_remaining()}.
  \item[$\square$] Overall budget $\geq$ attempts $\times$ attempt p99 + backoff.
  \item[$\square$] Handlers re-check budget after any stall.
  \item[$\square$] Stream loops poll \texttt{context.is\_active()}.
  \item[$\square$] Clients \texttt{cancel()} abandoned streams explicitly.
  \item[$\square$] Retry policy per method; UNAVAILABLE only; idempotency known.
\end{itemize}

\subsection*{Wire Cheat Sheet}
\begin{minted}{text}
message  = [compressed:1B][length:4B big-endian][protobuf payload]
:path    = /package.Service/Method     content-type = application/grpc
status   = grpc-status trailer (HTTP :status is 200 even on errors)
details  = <key>-bin trailer (raw bytes; base64 on the HTTP/2 wire)
deadline = grpc-timeout header, remaining budget, unit suffix H/M/S/m/u/n
\end{minted}

\section{Standardized Evidence Addendum}
\subsection{Benchmark Snapshot Template}
\begin{table}[h]
  \centering
  \small
  \begin{tabularx}{\textwidth}{l c c c c}
    \toprule
    \textbf{Config} & \textbf{p50 latency} & \textbf{p99 latency} & \textbf{Error rate} & \textbf{Throughput} \\
    \midrule
    Baseline & -- & -- & -- & 1.00x \\
    Candidate A & -- & -- & -- & -- \\
    Candidate B & -- & -- & -- & -- \\
    \bottomrule
  \end{tabularx}
  \caption{Minimum benchmark table to keep per chapter.}
\end{table}

\subsection{Request Trace Template}
\begin{itemize}
  \item \textbf{Request:} one representative API call for this chapter's topic.
  \item \textbf{Before:} behavior (headers, payload, latency, failure mode) with the naive approach.
  \item \textbf{After:} behavior with the technique in this chapter.
  \item \textbf{Result:} what changed in correctness, resilience, latency, and operability.
\end{itemize}

\subsection{Cost and Latency Budget Snapshot}
\begin{table}[h]
  \centering
  \small
  \begin{tabularx}{\textwidth}{l c c}
    \toprule
    \textbf{Stage} & \textbf{Latency budget} & \textbf{Cost driver} \\
    \midrule
    TLS / connection setup & -- & handshake round trips, keep-alive hit rate \\
    Request validation & -- & schema complexity, CPU per request \\
    Business logic and I/O & -- & downstream fan-out, query cost \\
    Serialization and egress & -- & payload size, compression ratio \\
    \bottomrule
  \end{tabularx}
  \caption{Operational budget template for this chapter's technique.}
\end{table}

\section{Configuration Preset Template}
\begin{minted}{text}
# api_policy.yaml
policy_version: "v1.0.0"
component: "<chapter_topic>"
timeout_ms: 0
retry_budget: 0
fallback_strategy: "fail_fast"
rollback_trigger: "error_budget_burn_gt_threshold"
\end{minted}

\section{CI Quality Gate Specification}
\begin{itemize}
  \item contract tests green against the pinned OpenAPI/AsyncAPI/Protobuf spec,
  \item no breaking schema changes without a version bump,
  \item error responses conform to RFC 9457 problem-details shape,
  \item benchmark deltas within approved regression bounds (latency and throughput),
  \item policy version and rollback plan present in release metadata.
\end{itemize}

\section{Incident Runbook Template}
\begin{enumerate}
  \item Detect regression from SLO burn-rate alerts or consumer reports.
  \item Scope impacted endpoints, versions, and consumer segments.
  \item Contain impact (rollback, feature flag, rate-limit tightening, or traffic shift).
  \item Diagnose with traces, structured logs, and contract diffs.
  \item Recover with validated fix and verified rollout procedure.
  \item Prevent recurrence via new regression tests and CI gates.
\end{enumerate}

\section{Cross-Chapter Decision Card}
\begin{table}[h]
  \centering
  \small
  \begin{tabularx}{\textwidth}{l X X}
    \toprule
    \textbf{Dimension} & \textbf{Choose this approach when} & \textbf{Prefer an alternative when} \\
    \midrule
    Use case & -- & -- \\
    Latency budget & -- & -- \\
    Operational cost & -- & -- \\
    Team maturity & -- & -- \\
    \bottomrule
  \end{tabularx}
  \caption{Decision card to complete per chapter.}
\end{table}

\section{ROI Model and Build-vs-Buy}
\begin{itemize}
  \item \textbf{Build when:} the capability is differentiating, compliance forbids vendors, or scale makes vendor pricing irrational.
  \item \textbf{Buy when:} the capability is commodity (auth, gateways, observability), time-to-market dominates, or staffing is thin.
  \item \textbf{Hidden costs to model:} on-call load, upgrade cadence, expertise ramp, data egress fees.
\end{itemize}

\section{Disaster Recovery Notes (RPO/RTO)}
\begin{itemize}
  \item Define recovery point objective (RPO): maximum acceptable data loss window.
  \item Define recovery time objective (RTO): maximum acceptable restoration time.
  \item Test restores regularly; an untested backup is not a backup.
\end{itemize}

